\documentclass[11pt,letterpaper]{article}
\usepackage[T1]{fontenc}
\usepackage[utf8]{inputenc}
\usepackage{lmodern,microtype}
\usepackage[margin=0.88in,headheight=15pt]{geometry}
\usepackage{amsmath,amssymb,amsthm,mathtools}
\usepackage{booktabs,longtable,tabularx,array,enumitem,fancyhdr,xcolor,listings,needspace}
\usepackage[hyphens]{url}
\usepackage{hyperref,bookmark}
\definecolor{ink}{RGB}{28,47,62}
\definecolor{pale}{RGB}{247,249,250}
\hypersetup{colorlinks=true,linkcolor=ink,citecolor=ink,urlcolor=blue!50!black,
  pdftitle={A Smaller Checked Presentation of Intuitionistic Primitive Recursive Arithmetic},
  pdfauthor={Research prepared for Vladimir Reshetnikov}}
\setlength{\emergencystretch}{3em}
\setlength{\parskip}{0.25em}
\setlist{itemsep=0.2em,topsep=0.35em,leftmargin=1.6em}
\newtheorem{theorem}{Theorem}[section]
\newtheorem{lemma}[theorem]{Lemma}
\newtheorem{proposition}[theorem]{Proposition}
\newtheorem{corollary}[theorem]{Corollary}
\theoremstyle{definition}\newtheorem{definition}[theorem]{Definition}
\theoremstyle{remark}\newtheorem{remark}[theorem]{Remark}
\newcommand{\ipra}{\mathsf{iPRA}}
\newcommand{\code}[1]{{\small\nolinkurl{#1}}}
\newcommand{\Rcal}{\mathcal R}
\newcommand{\Ncal}{\mathcal N}
\lstset{basicstyle=\ttfamily\scriptsize,columns=fullflexible,keepspaces=true,
  breaklines=true,showstringspaces=false,frame=single,rulecolor=\color{black!15},
  backgroundcolor=\color{pale},xleftmargin=0.3em,xrightmargin=0.3em}
\pagestyle{fancy}\fancyhf{}
\fancyhead[L]{\small Concise iPRA: syntax and reversible elimination}
\fancyhead[R]{\small Pinned comparison and proofs}
\fancyfoot[C]{\thepage}
\title{\color{ink}\textbf{A Smaller Checked Presentation of\\Intuitionistic Primitive Recursive Arithmetic}\\[0.6em]
  \large Shared syntax validation and reversible elimination}
\author{Research prepared for Vladimir Reshetnikov}
\date{November 11, 2026}
\begin{document}
\maketitle
\begin{abstract}
We give a complete presentation with the same consequences as the current
ConciseIPRA native kernel. Reversing the root equation after elimination
makes unrestricted equality reasoning derivable from a single transparent
term wrapper. The new source-elimination proof uses neither arithmetic
induction nor successor computation. The established whole-domain,
actual-equality bi-interpretation with full first-order iPRA therefore
transfers unchanged, including both formula round trips.

The implementation improvement is separate and explicitly measured. A
single parity-indexed syntax validator replaces mutually recursive term
and formula validators. The recommended standalone checker, including
exact integer/list validation of the entire certificate and validation
of the closed public goal, has 1,055 normalized bytes, versus 1,365 for a
matched control. Reversal itself costs six bytes. Under the repository's
more permissive compact convention, a neutral-premise change improves
526 to 525 bytes and, in a separate program, 283 to 282 AST nodes; the
259-token frontier is not improved. These are modest, reproducible
improvements, not a global minimality theorem. The new mathematics is
proved in writing; the supplied Lean extension has not been compiled.
\end{abstract}

\section{Result, benchmark, and attribution}
\label{sec:result}

The supplied repository was inspected at commit
\begin{center}\code{1517d9d79ce33f631e0f82960c849427cbba23bf}.\end{center}
All comparisons in this report are to that frozen snapshot, not to an
earlier draft or to a moving branch. Its README describes the native
kernel, twenty-eight research packages, and the existing Lean
bi-interpretation~\cite{baseline}. The smallest listed compact coordinates
are 259 significant tokens, 283 AST nodes, and 526 normalized bytes,
attained by different programs in \code{source_dispatch}~\cite{source}.

There are three distinct contributions.
\begin{enumerate}
\item \textbf{An elementary equality normalization theorem.} Change
  transport and closure elimination to reverse a root equation in their
  conclusion. A constant template now reverses any equation immediately.
  This removes the arithmetic construction needed to repair a literal-zero
  source in the older, unmodified transport rule. The whole new calculus
  is contextually equivalent to the current kernel.
\item \textbf{A smaller complete goal validator.} Term and formula
  validation share one recursion, with a depth invariant encoding their
  sorts. This is an implementation simplification available to either
  calculus. Matched controls quantify its contribution independently of
  reversal.
\item \textbf{A small compact-checker reduction.} Use scalar formula
  zero, an administrative truth, as the neutral retained premise. This
  changes one representation choice and saves one normalized byte, one
  token, and one AST node in each applicable local-scope profile. It does
  not change any arithmetic or logical axiom.
\end{enumerate}

The numerical language, the two computation equations, the conditional
binders, the reconstruction of intuitionistic logic, predecessor,
equality reflection, pairing, and the primitive-recursion compiler are
inherited. They are attributed and reproduced in the appendices. The
report does not claim to have rediscovered these results.

\paragraph{What ``simpler'' means here.}
The recommended program is smaller than a matched program with the same
input checks. The new equality argument has fewer mathematical
dependencies. Neither statement says that the new reversal operation
reduces the inherited compact kernel: it adds six normalized bytes.
The one-byte compact improvement belongs to the equivalent unreversed
calculus. No single delivered program attains all coordinate minima,
and none improves the current 259-token record. Absolute simplestness
would require a fixed description language and a lower-bound theorem;
the attained source sizes are upper bounds, not lower bounds.

\paragraph{Reproducibility and scope.}
The archive contains the report source and PDF, standalone checkers,
comparison controls, a certificate producer and fixtures, a reproduction
script, the pinned Lean sources, and an explicitly uncompiled Lean
extension. The source inventory verifies Git blob identities for the
frozen files it lists. A novelty screen inspected the README and the
eighty small primary Python modules retained in
\code{reference/current_small_checkers/}; it is not a literature-wide
priority search.

\clearpage
\setcounter{tocdepth}{1}
\begingroup\small\setlength{\parskip}{0pt}\tableofcontents\endgroup
\clearpage

% Includable fragment. Requires amsmath, amssymb, amsthm and the
% theorem/lemma/proposition/corollary environments. Labels use rev:.
% Written proof; the accompanying suggested Lean module has NOT been compiled.

\section{Reversal makes the source restriction elementary}
\label{rev:section}

This section replaces the current kernel's equational transport by a
uniform crossed elimination. Its purpose is to simplify the equality
argument while preserving exactly the native consequence relation. The
proof uses no arithmetic induction: a single, syntactically nonzero
retraction context supplies everything needed to remove the restriction
on the source of an equality.

\subsection{Complete calculus and ordinary reading}

Numerical terms and native formulas are unchanged:
\[
 t::=x\mid0\mid S(t)\mid I(n,a;u.h),
 \qquad
 A::=(s=t)\mid Q_yA\mid D_{x,y}A.
\]
Only $u$ is bound in the body $h$ of an iterator. The binders $Q$ and $D$
bind respectively $y$ and $(x,y)$ in their matrices. All substitutions
below are capture-avoiding; unused template variables are permitted.
The ordinary readings are
\[
 Q_yB(y)\equiv B(0)\to\forall y\,B(y),\qquad
 D_{x,y}B(x,y)\equiv
 \exists x\bigl(B(x,0)\to\forall y\,B(x,y)\bigr).
\]
The designated contradiction remains $0=S0$.

Define an operation on the root of a formula by
\[
 (s=t)^\dagger=(t=s),\qquad
 (Q_yB)^\dagger=Q_yB,\qquad
 (D_{x,y}B)^\dagger=D_{x,y}B.
 \tag{Reversal}\label{rev:rho}
\]
It does not recurse into a compound matrix. In particular,
$(A^\dagger)^\dagger=A$, and reversal commutes with capture-avoiding
numerical substitution. Neither statement expresses an object-theory
axiom: these are identities of finite syntax.

Write $s\not\equiv0$ to mean that $s$ is not the literal zero term.
This condition does not assert $s\ne0$ in arithmetic. In particular,
$I(0,0;u.0)\not\equiv0$, although its value is zero.

The calculus $\mathcal R$ has exactly the following eight rule schemas.
The retained context $\Gamma$ is a finite reusable list of native
formulas.
\[
 \frac{}{\Gamma\vdash A}\ (A\in\Gamma),\qquad
 \frac{}{\Gamma\vdash I(0,a;u.h)=a},\qquad
 \frac{}{\Gamma\vdash I(Sn,a;u.h)=h[I(n,a;u.h)/u]}.
 \tag{Hyp, $I_0$, $I_S$}\label{rev:initial}
\]
\[
 \frac{\Gamma\vdash s=t\qquad\Gamma\vdash F(s)}
      {\Gamma\vdash F(t)^\dagger}
 \quad(s\not\equiv0),
 \tag{X}\label{rev:x}
\]
where $F$ is any native formula template, including a constant template,
and
\[
 \frac{\Gamma\vdash Q_yB(y)\qquad\Gamma\vdash B(0)}
      {\Gamma\vdash B(t)^\dagger}.
 \tag{QX}\label{rev:qx}
\]
The remaining three rules are exactly those of the current native kernel:
\[
 \frac{\Gamma\vdash D_{x,y}B(x,y)
       \qquad\Gamma,Q_yB(b,y)\vdash A}
      {\Gamma\vdash A},
 \tag{DE}\label{rev:de}
\]
\[
 \frac{\Gamma,B(t,0)\vdash B(t,Sz)}
      {\Gamma\vdash D_{x,y}B(x,y)},
 \tag{DI-case}\label{rev:di}
\]
\[
 \frac{\Gamma,E(z)\vdash E(Sz)}
      {\Gamma\vdash Q_yE(y)}
 \quad(E\text{ a literal equation}).
 \tag{QI-progress}\label{rev:qi}
\]
In DE, $b$ is fresh for $\Gamma$, $A$, and the fixed parameters of $B$.
In DI-case and QI-progress, $z$ is fresh for the retained context and
fixed parameters; in DI-case it is also fresh for the chosen witness
$t$. The displayed local assumption is discharged. No logical axiom,
reflexivity axiom, cut rule, unrestricted substitution rule, numerical
case rule, separation axiom, or additional induction rule is primitive.
Their admissibility, when needed, follows from the equivalence proved
below. In particular, no rule-by-rule substitution into a derivation of
$\mathcal R$ is assumed in this proof.

\subsection{Reversal and the equality algebra}

Throughout this subsection, all displayed derivations are over an
arbitrary fixed context $\Gamma$. Introduce the abbreviations
\[
             J(a):=I(0,a;u.0),\qquad c:=J(0).
 \label{rev:j}
\]
The equation $J(a)=a$ is an instance of $I_0$. Both $J(a)$ and $c$
are syntactically nonzero, independently of the value of $a$.

\begin{lemma}[Unrestricted reversal]\label{rev:reverse}
From $\Gamma\vdash A$ one can derive $\Gamma\vdash A^\dagger$,
for every native formula $A$.
\end{lemma}
\begin{proof}
Apply X with major premise $c=0$, supplied by $I_0$, and the constant
template $F(v):=A$. Its minor premise is the given $A$, and its
conclusion is $A^\dagger$. The source condition holds because
$c\not\equiv0$. This use of a constant template is essential.
\end{proof}

\begin{corollary}[Uncrossed rules]\label{rev:uncrossed}
Ordinary QE is admissible. Ordinary transport
$s=t,F(s)\Rightarrow F(t)$ is admissible whenever $s\not\equiv0$.
\end{corollary}
\begin{proof}
Apply QX or X, respectively, and then Lemma~\ref{rev:reverse}.
The two reversals cancel syntactically.
\end{proof}

\begin{lemma}[Reflexivity and symmetry]\label{rev:refl-symm}
Every $a=a$ is derivable. From $a=b$ one can derive $b=a$, with no
restriction on either endpoint.
\end{lemma}
\begin{proof}
For reflexivity use $J(a)=a$ twice, as the major and minor premises of X,
with template $F(v):=(v=a)$. The conclusion is $(a=a)^\dagger=a=a$.
Symmetry is the equational case of Lemma~\ref{rev:reverse}. Notice that
this includes equations with the literal source $0$.
\end{proof}

\begin{lemma}[Joining at zero]\label{rev:zero-join}
If $p,q\not\equiv0$, then from $p=0$ and $q=0$ one can derive $p=q$.
\end{lemma}
\begin{proof}
Use major premise $p=0$, minor premise $J(p)=p$, and template
$F(v):=(J(v)=p)$ in X. Its conclusion is
\[
                     (J(0)=p)^\dagger=(p=c).
\]
Similarly obtain $q=c$. Reverse both equations to get $c=p$ and
$c=q$. Finally apply X with major $c=q$, minor $c=p$, and template
$F(v):=(v=p)$. Its conclusion is $(q=p)^\dagger=(p=q)$.
Every crossed source used here is one of $p,q,c$, all syntactically
nonzero. No substitution into a proof, induction, or numerical case
distinction occurs.
\end{proof}

\begin{lemma}[Unrestricted transitivity]\label{rev:trans}
From $r=s$ and $s=t$ one can derive $r=t$.
\end{lemma}
\begin{proof}
If $s\not\equiv0$, apply the uncrossed restricted transport of
Corollary~\ref{rev:uncrossed} to major $s=t$, minor $r=s$, and
template $F(v):=(r=v)$.

Suppose instead that $s\equiv0$. If $r\equiv0$, the second premise is
already the conclusion. If $t\equiv0$, the first premise is already
the conclusion. In the remaining case $r,t\not\equiv0$, the premises
are $r=0$ and $0=t$. Reverse the latter and apply
Lemma~\ref{rev:zero-join}.

The distinction between a literal zero and another term is a decidable
case distinction on finite syntax in the metaproof. It is not an
application of excluded middle or induction in the object calculus.
\end{proof}

\begin{lemma}[Unrestricted numerical congruence]\label{rev:congr}
For every numerical template $f(v)$, from $s=t$ one can derive
$f(s)=f(t)$.
\end{lemma}
\begin{proof}
If $s\not\equiv0$, apply uncrossed restricted transport to major $s=t$
and template $F(v):=(f(s)=f(v))$, starting from reflexivity
$f(s)=f(s)$.

If $s\equiv0$ and $t\not\equiv0$, reverse the major to $t=0$ and use
template $F(v):=(f(v)=f(t))$. Reflexivity $f(t)=f(t)$ is the required
minor premise, so the conclusion is $f(0)=f(t)$.

If both endpoints are literally zero, the required equation is an
instance of reflexivity.
\end{proof}

\begin{theorem}[Recovery of equational transport]\label{rev:tr-eq}
The rule
\[
 \frac{\Gamma\vdash s=t\qquad\Gamma\vdash a(s)=b(s)}
      {\Gamma\vdash a(t)=b(t)}
\]
is admissible in $\mathcal R$, without a restriction on $s$.
\end{theorem}
\begin{proof}
Numerical congruence gives $a(s)=a(t)$ and $b(s)=b(t)$. Reverse the
first equation and use unrestricted transitivity twice:
\[
                         a(t)=a(s)=b(s)=b(t).
\]
\end{proof}

\subsection{The abstract retraction argument}

The preceding proof has a stronger, arithmetic-independent form. This
form identifies precisely which facts are used, and is also a concise
way to check that the equality argument is not circular.

\begin{proposition}[Retraction criterion]\label{rev:abstract}
Suppose a term language has a distinguished literal $0$ and a fixed
one-variable term context $J(v)$ satisfying both of the following
conditions:
\begin{enumerate}
\item $J(a)$ is syntactically different from $0$ for every term $a$;
\item $J(a)=a$ is an available axiom schema, uniformly in the retained
context.
\end{enumerate}
For equations, let reversal exchange their endpoints. Then X with
sources restricted to terms other than literal $0$ derives all ordinary
equational transport. The derivation uses only the retraction schema
and X.
\end{proposition}
\begin{proof}
Set $c=J(0)$. A single use of X with constant template $A$ and major
$c=0$ reverses any equation $A$. The retraction instance $J(a)=a$ used
twice derives reflexivity. For a nonzero term $p$ satisfying $p=0$,
the template $J(v)=p$ turns its retraction instance into $p=c$.
Thus any two nonzero terms equal to $0$ can be joined through the
nonzero anchor $c$. Ordinary transitivity now follows by splitting on
whether its middle term is literally $0$. Congruence follows by using
the nonzero endpoint of the source equation, reversing its orientation
when necessary; the equation $0=0$ is handled by reflexivity. Combining
congruence, symmetry, and transitivity gives general equational
transport, as above.
\end{proof}

For $\mathcal R$, take $J(v)=I(0,v;u.0)$. Consequently the entire
equality normalization theorem uses neither $I_S$, nor QX, nor DE,
nor either introduction rule. In particular, the numerical induction
rule is not invoked even indirectly. This contrasts with the supplied
ordinary source-restricted transport proof, whose zero-source symmetry
construction uses numerical tests and equation induction. The present
argument applies to a changed crossed rule, rather than giving a
shorter proof of that older rule's difficult intermediate lemma.

\subsection{Exact equivalence with the current kernel}

Let $\mathcal N$ denote the native kernel formalized as
\texttt{NPrf} at the pinned source commit. Its equality rule is
equational transport, its closure elimination is ordinary QE, and its
remaining six rules are Hyp, $I_0$, $I_S$, DE, DI-case, and QI-progress.

\begin{theorem}[Identity-language consequence equivalence]
\label{rev:equivalence}
For every finite native context $\Gamma$ and formula $A$,
\[
                 \Gamma\vdash_{\mathcal R}A
                 \quad\Longleftrightarrow\quad
                 \Gamma\vdash_{\mathcal N}A.
\]
There are effective, syntax-preserving translations of finite
derivations in both directions.
\end{theorem}
\begin{proof}
For $\mathcal R\to\mathcal N$, translate each finite derivation by
induction on its construction. The six unchanged rules translate
directly. For X, use the existing derived theorem
\texttt{NPrf.tr}, which transports through an arbitrary native
template, and then reverse the resulting equality if the formula is
atomic. If its root is $Q$ or $D$, reversal is the identity. For QX,
use ordinary QE followed by the same reversal. Equality symmetry is an
existing derived rule of $\mathcal N$.

For $\mathcal N\to\mathcal R$, again recurse on the derivation. The
six common rules are unchanged, equational transport is supplied by
Theorem~\ref{rev:tr-eq}, and ordinary QE is supplied by
Corollary~\ref{rev:uncrossed}. These constructions work over arbitrary
retained contexts, so they apply unchanged inside the fresh-variable
subderivations of DE and the two introduction rules.

There is no need to reprove compound transport inside $\mathcal R$:
the old kernel requires only its equational case in this direction.
The existing compound-transport theorem is used only for the
$\mathcal R\to\mathcal N$ translation, where it is already available.
\end{proof}

\begin{corollary}[Bi-interpretability and equiconsistency]
\label{rev:biinterpretation}
The existing native interpretations into and out of full first-order
iPRA are also a bi-interpretation between $\mathcal R$ and iPRA.
The numerical domains remain whole, equality retains its ordinary
meaning, and both interpretation composites have the same identity
graph $x=y$ as before. In particular, $\mathcal R$ is equiconsistent
with iPRA.
\end{corollary}
\begin{proof}
The term and formula languages, ordinary readings, and designated
contradiction have not changed. Apply
Theorem~\ref{rev:equivalence} to every sequent used in the existing
native-to-iPRA faithfulness and round-trip theorems. This transfers the
theorems \texttt{native\_iff\_pra}, \texttt{pra\_iff\_native},
\texttt{native\_roundTrip}, \texttt{pra\_native\_roundTrip}, and
\texttt{native\_inconsistent\_iff}; no new numerical representation
or quotient is introduced.
\end{proof}

\paragraph{Verification boundary.}
The argument above is a complete written derivation relative to the
supplied native kernel and its already established iPRA bridge. The
accompanying \texttt{Reversal.lean} is a proposed formalization using
the current source names and deeply embedded rules. It contains no
\texttt{sorry} or added axiom, but it has not been compiled in this
environment and is not claimed as machine verification. Checker-size
measurements are separate from this mathematical equivalence theorem.

\section{The interpretations, including both composites}
\label{sec:interpretations}

The equivalence in Theorem~\ref{rev:equivalence} is stronger than agreement
on standard numerical computations: it holds for every sequent, with
arbitrary retained assumptions. This is what permits reuse of the
existing bi-interpretation rather than only an equiconsistency argument.
For clarity we spell out the translations and the role of the inherited
arithmetic results.

\subsection{Which iPRA is intended}

Here iPRA is full first-order intuitionistic logic with equality, a
function symbol for every fixed primitive-recursive definition, its
defining equations, the successor axioms, and induction for
quantifier-free formulas. Quantified formulas are allowed everywhere
in the logic; induction is not extended to arbitrary quantified formulas.
The pinned \code{PR.lean} implements this convention~\cite{leanbridge}.

The binding notation $I(n,a;u.h)$ denotes a fixed finite description.
It can be elaborated into a staged first-order signature: for each
earlier finite body $h(u,\vec p)$ introduce $i_h(n,a,\vec p)$ with
the two iterator equations. Parameters in the description become
ordinary first-order arguments. No variable ranges over functions or
programs, and no uniformly total evaluator for all primitive-recursive
descriptions is asserted.

\subsection{Native formulas and ordinary intuitionistic logic}

Let $T_I$ be ordinary intuitionistic first-order logic over the scalar
terms, with the two computation equations, induction for literal
equations, coverage
\[
  \forall x\,(x=0\vee\exists y\,x=Sy),
\]
and the explicit separation bridge $(0=S0)\to\bot$. Translate native
equations identically, and translate $Q$ and $D$ by the meanings given
in Section~\ref{rev:section}. Both interpretations use the full domain,
defined by $x=x$, and interpret equality as equality.

The reverse logical translation uses only native formulas. Put
$C(v,a,b)=I(v,a;u.b)$, with the state variable unused. Its values at zero
and successors are $a$ and $b$. The gate $G_vA$ is true at zero and
equivalent to $A$ at a successor; the selector $M_v(A,B)$ selects $A$
at zero and $B$ at successors. Both are constructed by external
recursion on fixed finite formulas, using equation selectors and
padding to align binder depth. Their full definitions and proof rules
are reproduced in Appendix~\ref{sec:choice-introduction}.

With $p$ the parity iterator, one convenient Eq/$D$ macro inventory is
\[
\begin{aligned}
 \bot_*&=(0=S0),\\
 A\to_*B&=D_{u,z}M_z(A,B),\\
 A\land_*B&=D_{u,z}G_zM_{p(z)}(A,B),\\
 \exists_*x A(x)&=D_{x,y}G_yA(x),\\
 A\lor_*B&=\exists_*x M_x(A,B),\\
 \forall_*x A(x)&=A(0)\land_*D_{u,z}G_zA(z).
\end{aligned}
\]
Here the unused coordinate $u$ and all new bound variables are chosen
fresh. These are definitions, not new axiom schemas. The logical
introduction and elimination rules, falsity elimination, numerical
coverage, and both native/ordinary formula composites follow by the
explicit constructions in the appendix. An auxiliary $Q_yB$ is
interderivable with $D_{u,y}B$ when $u$ is unused.

The ordinary separation bridge must not be omitted: positive numerical
equations alone admit a one-element model. Native falsity is a designated
equation, and the bridge identifies its intended logical role in $T_I$.
This is already part of the inherited presentation, not an additional
assumption introduced for reversal.

\subsection{Scalar arithmetic and primitive recursion}

Interpret an iterator in iPRA by primitive recursion whose update ignores
the recursion counter. This translates the scalar equations and literal
equation induction directly. Coverage and separation are iPRA theorems.

For the converse translation, the inherited internal derivations proceed
in the following dependency order.
\begin{enumerate}
\item A scalar predecessor $P$ is constructed and proved to satisfy
  $P0=0$ and $P(Sn)=n$ using equation induction and numerical cases.
  Pairing and successor injectivity are not assumed in that bootstrap.
\item Truncated subtraction and a Boolean reflection of actual equality
  are derived. Each fixed quantifier-free formula is represented by a
  characteristic equation, giving its induction instance.
\item With quantifier-free induction available, pairing $\pi$ and
  projections $L,H$ are constructed internally. They satisfy
  $L(\pi(a,b))=a$ and $H(\pi(a,b))=b$ throughout the domain.
\item Each fixed primitive recursion is compiled by storing its counter
  and current value in one paired scalar state.
\end{enumerate}

For example, for
\[
 f(0,\vec p)=g(\vec p),\qquad
 f(Sn,\vec p)=h(n,f(n,\vec p),\vec p),
\]
the translation uses
\[
\begin{split}
 V(n,\vec p)&=I\bigl(n,\pi(0,g^*(\vec p));\,
 u.\pi(SL(u),h^*(L(u),H(u),\vec p))\bigr),\\
 f^*(n,\vec p)&=H(V(n,\vec p)).
\end{split}
\]
Equation induction proves $L(V(n,\vec p))=n$, and the projection
identities establish the defining equations of $f$. Translation is
external recursion on a fixed finite definition, not an internal
universal-recursion principle. The complete inherited constructions
appear in the scalar predecessor, reflection, and primitive-recursion
appendices.

\subsection{Why both round trips are identities}

For a fixed iterator, equation induction proves uniqueness from its
zero and successor equations. The analogous uniqueness proof applies
to each fixed primitive recursion. Inducting externally over a finite
function description therefore proves that its double translation is
equal to the original function inside the corresponding theory. The
interpretation domain is never changed, so the graph of the resulting
domain isomorphism is literally $x=y$.

For formulas, use the proved meanings of the logical macros and the
native $Q,D$ binders, then induct on formula structure. For example,
the native paired composite is
\[
 D_{x,y}B\ \dashv\vdash\
 \exists_*x\bigl(B(x,0)\to_*\forall_*y B(x,y)\bigr).
\]
The forward direction opens a fresh witness and introduces the three
logical macros; the backward direction eliminates the existential
macro and uses DI-case with that witness. The appendix gives both
directions with their local assumptions.

Combining these arithmetic and logical composites with
Theorem~\ref{rev:equivalence} establishes the requested bi-interpretation
for $\Rcal$. In particular,
\[
 \varnothing\vdash_{\Rcal}0=S0
 \quad\Longleftrightarrow\quad
 \ipra\vdash\bot.
\]
This is a relative consistency equivalence; no assertion of the
consistency of either theory is needed as a premise.

\subsection{Direct closure versus the implemented progress packet}
\label{sec:progress-packet}

The mathematical presentation in Section~\ref{rev:section} follows the
current Lean kernel and introduces $Q_yE(y)$ directly by equation
progress. The tiny inherited certificate format instead introduces
$D_{x,y}E(x,y)$ from $E(t,z)\vdash E(t,Sz)$, with a witness $t$.
These are equivalent interfaces, not literally the same instruction.

To obtain witnessed progress from direct closure, set $F(y)=E(t,y)$.
QI-progress gives $Q_yF(y)$. Under $F(0)$, ordinary QE gives $F(Sz)$;
DI-case with witness $t$ finishes. Ordinary QE is already derived in
$\Rcal$ by reversing QX's result.

Conversely, apply witnessed progress with an unused witness coordinate
$u$ to obtain $D_{u,y}E(y)$. Eliminate it. The supplied assumption is
exactly $Q_yE(y)$, independent of the fresh witness, so hypothesis
closes the branch. Both constructions are valid in arbitrary retained
contexts and keep the progress matrix a literal equation. The supplied
producer uses the witnessed packet convention; the Lean draft uses the
direct-closure convention.

\section{A smaller complete checker and its exact contract}
\label{sec:checker}

\subsection{Three input contracts must not be conflated}

All programs are whole standalone Python modules. They read a JSON
array, consume instruction packets from its right-hand end, and print
exactly \code{ACCEPT} on normal successful termination. Exceptions and
nonzero exit codes are rejection. Unused prefix packets are inert.
The recommended program additionally requires every node of the entire
JSON value, including an unused prefix, to be an exact integer or list.

\begin{center}\small
\begin{tabularx}{\linewidth}{@{}lX@{}}
\toprule
Contract & What the program checks\\
\midrule
Compact & The inherited proof protocol; canonical public statements and
the advertised certificate representation are a caller contract.\\
Goal-validating & The protocol plus an explicit canonical closed-goal
parser. This does not assert that every internal expression is a canonical
object formula: administrative aliases remain part of the protocol.\\
Fully tree-typed & The goal-validating contract plus exact integer/list
syntax throughout the whole JSON certificate. Booleans, floats, null,
strings, and objects are rejected before proof checking.\\
\bottomrule
\end{tabularx}
\end{center}

The recommended \code{code/checker_rev_typed.py} uses the last contract
and admits public Eq/$Q$/$D$ goals. Its rule certificates implement the
witnessed-progress version explained in Section~\ref{sec:progress-packet}.
The compact token/AST experiments are not substituted for this explicit
input policy.

\subsection{The arity representation}

The inherited representation shares lists between numerical expressions
and formulas. Their sort is determined by position. The table describes
the relevant canonical constructors; the exact slot encoder is included
with the certificate producer.
\begin{center}\small
\begin{tabular}{@{}lll@{}}
\toprule
Raw representation & Numerical position & Formula position\\
\midrule
$0$ & Numerical zero & Administrative truth\\
$[a]$ & Successor & $D$ binder\\
$[a,b]$ & Zero alias & Equation\\
$[a,b,c]$ & Scalar iterator & Closure resource\\
$[B,B,B]$ & --- & Canonical $Q_yB$\\
Empty or longer lists & Zero aliases where applicable & Truth padding\\
\bottomrule
\end{tabular}
\end{center}

Public formulas contain only the canonical Eq/$Q$/$D$ forms, and their
numerical leaves are zero or in-scope protected slot indices. The codec
inserts dummy zero slots so that the opening routine can treat all
children of a node uniformly. Protected indices are odd numbers starting
at 3; rigid local parameters use allocated negative integers. Dummy slots
denote zero, not extra free variables. The proof checker validates images
against the current local freshness floor before substitution.

An instruction is a quadruple $[m,s,u,t]$. The major $m$ is checked
recursively before any supplied resource is trusted. Source $s$ selects
the expected elimination major:
\[
 M(s,u,t)=
 \begin{cases}
 [s,u],&s\text{ is nonzero in the canonical raw representation},\\
 [t,t,t],&s=0.
 \end{cases}
\]
Thus a matched nonzero source performs X and a zero source performs QX.
For a matched elimination, the goal is the reversal of the target
instance and the continuation goal is the source instance. A proved
unary major exports its closure resource; other major arities export
only harmless padding. A neutral major instead supplies a computation
equation or opens a $D$ introduction, depending on the initializer/mode
field. These are the existing source-dispatch packet cases with the
displayed reversal and neutral-premise changes.

\subsection{The new shared syntax validator}

Let $F(a,s)$ and $T(a,s)$ be the separate formula and term validators in
the frozen strict Eq/$D$ control. The new validator $V(a,d)$ satisfies
\[
 F(a,s)\iff V(a,2s),\qquad T(a,s)\iff V(a,3+2s).
 \tag{Scope invariant}\label{eq:validator-invariant}
\]
An odd depth is numerical; an even depth is a formula depth. Recursing
through an array $a$ changes the depth by $5-|a|$, exactly as in the
inherited opener. The principal checks, in readable pseudocode, are
\begin{lstlisting}[language=Python]
if type(tree) is int:
    return depth % 2 and (tree == 0 or
                          3 <= tree < depth and tree % 2)
require type(tree) is list
require arity in (1, 3) at numerical depth
require arity in (1, 2, 3) at formula depth
require a formula triple to have three identical children
return all(valid(child, depth + 5 - arity) for child in tree)
\end{lstlisting}

\begin{proposition}[Validator equivalence]\label{prop:validator}
For every JSON value and every nonnegative scope count, the shared
validator agrees with the separate validators. This holds both for
the frozen Eq/$D$ grammar and for its explicitly constructed Eq/$Q$/$D$
extension.
\end{proposition}
\begin{proof}
Prove the two statements in \eqref{eq:validator-invariant} simultaneously
by structural induction. Exact type tests reject the same noninteger,
nonlist values. At numerical depth $3+2s$, both admit exactly zero and
odd indices $3\le a<3+2s$. At formula depth, both reject scalar leaves.
A unary formula adds four to the depth and binds two slots; binary
equality adds three and enters the numerical sort. Unary successor adds
four, accounting for its two dummy slots; ternary iteration adds two,
accounting for one protected slot. No other arity is accepted. For the
full grammar, a diagonal formula triple adds two, corresponding to one
$Q$ variable. The identical-child check is the same in both validators.
\end{proof}

The certificate-type predicate is separate and elementary: an integer
passes, and a list passes exactly when all its children pass. It is
counted in full in the recommended program and in its matched control.

\subsection{Scalar zero as a neutral formula}

At formula depth a scalar denotes administrative truth. The old local
checker begins with the retained neutral formula \code{[]}; the new
one begins with \code{0}. The new source is \code{g=[0]} instead of
\code{g=[[]]}, or \code{g=(0,)} instead of \code{g=([],)} in the tuple
profile. Each change saves one normalized character, token, and AST
node in the applicable profile.

This is sound because only the spelling of a true retained resource
changes. For completeness, emit neutral packet majors as zero; the
computation and introduction branches still see a falsy major. Their
case/progress selection is carried by the unchanged mode field. Native
public formulas do not contain an administrative truth leaf in formula
position. Consequently the represented native theorems are unchanged.
The serialized-certificate languages are not identical, and the
provided certificate conversion changes neutral packet fields explicitly.
Scalar formula zero is not the contradiction, which is the equation
$0=S0$, encoded as a binary formula.

\subsection{Why raw root reversal is safe}
\label{sec:raw-reversal}

For canonical formulas, array reversal swaps the two sides of an
equation and leaves a unary $D$ or a diagonal $Q$ triple unchanged.
Intermediate administrative triples need not be diagonal. It would be
incorrect to reverse such a triple while retaining the old convention
that only its last child supplies its meaning.

We use a symmetric administrative reading instead:
\[
 [B_0,B_1,B_2]\quad\rightsquigarrow\quad
 Q_yB_0\ \land\ Q_yB_1\ \land\ Q_yB_2.
 \tag{Triple reading}\label{eq:triple-reading}
\]
This is a conjunction of closures, not the closure of a conjunction.
A diagonal triple retains its ordinary $Q$ meaning by conjunction
idempotence. Opening acts on all three children at the same depth, so
substitution acts componentwise. Root reversal merely reverses the
conjunct order and is an intuitionistic equivalence.

All instruction cases respect this reading. The QX matcher requires an
exact diagonal triple. A unary $D$ major exports precisely such a
diagonal. A nonempty major of another array arity exports a list of
arity at least six, which is administrative truth. X is sound in any
decoded formula by equality and reversal. DI-case is sound for every
decoded matrix by numerical coverage; its arithmetic progress alternative
applies to genuine binary equations, or to harmless truth padding.
The arithmetic codec and its opening lemma are unchanged.

Thus arbitrary integer/list certificate aliases do not turn a canonical
false goal into a theorem by swapping a noncanonical triple. The changed
administrative reading is a real proof obligation, and is stated here
rather than concealed behind the public-goal parser. A larger alternative
can reverse only even roots and keep the old reading; it is not needed
for the delivered recommendation.

\subsection{Freshness and successful-trace invariant}

The retained context is finite and reusable. Each recursive branch has
its own floor $k$; allocated rigid parameters lie at or above that floor.
Before a fresh parameter is used, the floor is decreased, and only then
are the new assumption and successor target opened. The major premise
is checked at the old floor. This order prevents a future witness from
being captured by a parameter that has not yet been introduced.

For a successful canonical trace, a retained resource is either a decoded
assumption of the corresponding native derivation or a truth/computation
resource justified by the inherited supply rule. Each non-hypothesis
packet first verifies its major. A fresh auxiliary computation resource
whose zero instance is a theorem is admissible: substitute zero for the
fresh parameter in its continuation, then cut that theorem. The new
equality equivalence transfers substitution and cut from the native
kernel, so this argument does not assume a primitive substitution rule
for the restricted-source calculus.

A proposed shortcut replacing the mutable local floor by a static
context-length expression was rejected during this work. It did not
preserve the before/after opening phase and either rejected its own
fresh parameter or admitted it prematurely. Its attractive source-size
figures are not included as improvements and its code is not in the
deliverable.

\section{Whole-program measurements and matched controls}
\label{sec:measurements}

The measurements use CPython 3.12.14, matching the pinned repository's
convention. Remove genuine docstrings, apply \code{ast.unparse}, and add
one final newline. Count significant tokens, AST nodes, and UTF-8 bytes
of that entire normalized module. Standard-library JSON and the Python
runtime are common infrastructure; input checks, proof rules, freshness,
substitution, and the verdict driver are all included. Gzip uses maximum
compression, level 9, with \code{mtime=0} and no embedded filename. Its
size is reported for information and was not the optimization objective.

\begin{center}\small
\begin{tabular}{@{}lrrrr@{}}
\toprule
Complete program & Tokens & AST & Bytes & Gzip-9\\
\midrule
Old control, full typed contract & 532 & 618 & 1365 & 614\\
Shared parser and neutral zero, no REV & 450 & 514 & 1049 & 540\\
\textbf{Recommended typed REV} & \textbf{456} & \textbf{520} & \textbf{1055} & \textbf{545}\\
\midrule
Frozen old Eq/$D$ goal validator & 450 & 521 & 1109 & 538\\
Shared Eq/$D$ goal validator, no REV & 373 & 426 & 824 & 460\\
REV with Eq/$D$ goal validator & 379 & 432 & 830 & 466\\
Old full-goal control & 481 & 560 & 1189 & 554\\
Shared full-goal control, no REV & 399 & 456 & 873 & 477\\
REV with full-goal validator & 405 & 462 & 879 & 482\\
\bottomrule
\end{tabular}
\end{center}

The full-goal and full-typed old controls are newly constructed comparison
derivatives: they extend the frozen separate parser to canonical $Q$
triples and, for the typed control, add the same exact integer/list
predicate. They are not misrepresented as upstream files. The frozen
Eq/$D$ control and three source-dispatch controls are checked against
their upstream Git blob identifiers.

The recommended checker is 310 normalized bytes smaller than its
matched control, a reduction of approximately $22.71\%$. The
same-refactor control is six bytes smaller than the REV checker:
reversal is not responsible for the source-size gain. The Eq/$D$
comparison against an unchanged upstream whole program is also given
so that the conclusion does not depend solely on a newly extended
control.

\begin{center}\small
\begin{tabular}{@{}lrrrr@{}}
\toprule
Compact protocol, no public-goal validator & Tokens & AST & Bytes & Gzip-9\\
\midrule
Frozen byte profile & 266 & 305 & 526 & 330\\
Neutral-zero byte profile & 265 & 304 & \textbf{525} & 329\\
Frozen AST profile & 263 & 283 & 607 & 348\\
Neutral-zero AST profile & 262 & \textbf{282} & 606 & 348\\
Frozen token profile & \textbf{259} & 289 & 572 & 339\\
REV compact, neutral zero & 271 & 310 & 531 & 334\\
\bottomrule
\end{tabular}
\end{center}

The compact byte and AST improvements are one-coordinate changes to
existing profiles. They do not turn the two records into one program,
and they do not improve the token record. The 1,055-byte typed checker
and 525-byte compact checker have different input contracts and must
not be compared as if they did the same validation. The unreversed
compact checkers prove exactly the same native sequents as $\Rcal$
through the proved equivalence, but their instruction certificates use
the older elimination orientation.

\section{Executable evidence and verification limits}
\label{sec:validation}

The supplied finite validation covers actual derivations, public-goal
validation, and the administrative boundary. It is separate from the
universal written proofs.
\begin{itemize}
\item \textbf{REV certificates:} 36 cases, 362 process expectations across
  five profiles, ordinary Python and \texttt{python -O}, and a matched
  non-REV orientation control; zero mismatches.
\item \textbf{Readable checker:} another 72 process expectations on the
  same cases; zero mismatches. For the recommended typed interface each
  execution mode has 16 accepts and 20 rejects.
\item \textbf{Neutral premise:} 40 comparisons against frozen old
  certificates and checkers, in ordinary and optimized Python; zero
  mismatches.
\item \textbf{Validator equivalence:} 260,500 comparisons over generated
  canonical trees, malformed values, and five scope sizes; zero
  mismatches. The mathematical justification is the structural induction
  in Proposition~\ref{prop:validator}, not the sample count.
\end{itemize}

Positive certificates exercise the new reflexivity and constant-template
reversal, literal-zero-source symmetry, joining at zero, congruence
through nested iterator binders, unrestricted equation transport, QX
specialization, nested $Q$ cases, a nonvacuous $D$ witness, contextual
equation induction, and transport under both native binders. Negative
controls include a false goal, an unproved major, a removed induction
hypothesis, forbidden compound progress, unallocated witnesses and
specialization targets, free labels in closed goals, and malformed
public syntax. The exact REV orientation fixture is rejected by the
matched unreversed checker, so the test suite distinguishes the new
rule rather than exercising only common code paths.

\subsection{A concrete proof-size improvement}

The frozen \code{zero_symmetry_one_local.json} and the new
\code{conditional_zero_symmetry.json} have exactly the same raw closed
goal. Their compact certificate sizes are:
\begin{center}\small
\begin{tabular}{@{}lrr@{}}
\toprule
Certificate & Instruction packets & Compact JSON bytes\\
\midrule
Old zero-source normalization & 114 & 5435\\
REV construction & 37 & 1259\\
\bottomrule
\end{tabular}
\end{center}
This comparison demonstrates the practical effect of the elementary
reversal proof for one theorem. It is not a general proof-length bound
or a proof-checker-size measurement.

The mathematical proof uses $c=I(0,0;u.0)$ as its anchor. The packet
producer first uses $c_*=I(2,1;u.0)$, whose zero value follows from one
permitted successor packet with a nonzero initializer and predecessor.
It can then normalize computation leaves having a literal-zero
initializer or successor predecessor by replacing those positions with
$c_*$ and transporting along $c_*=0$. This prevents a circular call to
the very zero-initializer computation being normalized. This packet
normalization is inherited in method; the new equality derivations and
orientation compiler are explicit in \code{tests/rev_producer.py}.

\subsection{What has and has not been formally checked}

The pinned repository contains a Lean formalization of the inherited
kernel and its iPRA bi-interpretation. Its README reports successful
builds without placeholders~\cite{baseline,leanbridge}. The new
\code{formalization/Reversal.lean} is a proposed deeply embedded
extension using those exact source names. It contains proof bodies,
no added axioms, and no proof placeholders, but neither Lean nor Lake
was available in the authoring environment. It has not been compiled
and is not presented as machine verification. Its status is independent
of the completed Python executions and of the written theorem.

No external human peer review is claimed. The report establishes the
new mathematical result in writing relative to the pinned native
theory, includes the inherited bridge for inspection, and supplies
reproducible implementation evidence. It does not claim a new lower
bound on all possible proof checkers.

\section{An obstruction and a precise upstream wording correction}
\label{sec:limits}

One tempting deletion does not work: keep QI-progress for equations
but permit DI-case only at compound matrices, leaving the formula
grammar unchanged. This loses even $D_{x,y}(0=0)$.

To see this, interpret terms in standard $\mathbb N$, equations as
actual equality, and $Q$ by its ordinary reading. Interpret $D_{x,y}B$
as false when $B$ is literally an equation, and by its ordinary
existential conditional-closure reading otherwise. Numerical
substitution preserves the equation/compound distinction. Computation,
transport, hypotheses, QE, and equation progress are sound. DE is
vacuous at an atomic matrix and ordinary existential elimination
otherwise; the remaining DI-case instances are sound by coverage.
But $D_{x,y}(0=0)$ is false in this interpretation, although it is
provable in the original calculus. Thus that rule deletion cannot
replace the new construction. It would require a changed grammar or
an additional generating principle, as the existing guarded proposals
already recognize~\cite{guarded}.

A separate wording correction concerns the upstream restricted-source
calculus \code{ZPrf}. Substituting literal zero for a variable source
does not preserve a primitive restricted rule instance. Nevertheless
substitution is \emph{admissible} in its consequence relation, because
\code{zprf_iff_nprf} translates to the unrestricted kernel, where
substitution is admissible, and back. The failed property is direct
rule-by-rule preservation, not admissibility. This distinction does not
invalidate the existing equivalence theorem and is not used to shortcut
the new equality proof.

\section{Package and reproduction}

The archive is self-contained for the Python checks and report source.
The main \code{ipra_reversal_core.tex} embeds all sections and appendices.
To compile the Lean draft, copy it into the pinned repository's configured
Lean/Mathlib workspace; the included source snapshot does not replace
that dependency environment.

From the archive directory:
\begin{lstlisting}[language=bash]
python3 reproduce.py
python3 reproduce.py --pdf
python3 code/checker_rev_typed.py < tests/certificates/zero_neutral/conditional_zero_symmetry.json
\end{lstlisting}
The first command reproduces whole-program measurements and all
REV/readable and neutral-marker checks. The second also runs
\code{pdflatex} three times. The optional parser audit is
\texttt{python3 tools/audit\_validators.py}. Each checker is standalone and
imports only the standard library.

\clearpage
\appendix
\section{Attribution and reading the inherited appendices}

The following five mathematical sections are reproduced from the
pinned \code{certified_zero/reference/} sources under the repository's
MIT No Attribution license. Their historical statements such as
``the new rule'' or ``this paper'' refer to their original source.
They document the inherited Choice-DI interface and arithmetic bridge;
the present contribution is the equivalence in
Section~\ref{rev:section} and the implementation changes in
Section~\ref{sec:checker}. The correspondence between their witnessed
progress and the present direct-closure presentation is proved in
Section~\ref{sec:progress-packet}. No inherited proof is being claimed
as a new result of this report.

% New mathematical result: complementary introduction choices.
% Requires amsmath/amssymb/amsthm and theorem, lemma, proposition environments.
% The inherited arithmetic-to-iPRA bridge must be supplied separately.

\section{Complementary introductions with one local assumption}
\label{sec:choice-introduction}

The old paired introduction retains both the zero instance and an
induction instance when its matrix is an equation. Both together are
unnecessary. A certificate may choose either one, provided the induction
choice is restricted to a literal equation. This changes neither the
formula language nor the meaning of a formula. The omitted combination
is recovered by a short nested introduction--elimination construction.

\subsection{Complete native syntax and rule inventory}

The scalar numerical language is
\[
 t::=x\mid0\mid St\mid I(n,a;u.h).
\]
Only $h$ binds the state variable $u$; $n$ and $a$ do not.
The update may contain fixed numerical parameters. The computation
axioms in every retained context are
\[
 I(0,a;u.h)=a,\qquad
 I(Sn,a;u.h)=h[I(n,a;u.h)/u].
 \tag{I0, IS}\label{eq:choice-computations}
\]
The argument also applies unchanged to the full fixed-description PR
language with its projection, composition, and recursion equations.
No numerical variable ranges over function descriptions, and no uniform
evaluator is postulated.

Native formulas and reusable context clauses are
\[
 A::=(s=t)\mid\mathsf D_{x,y}A,\qquad
 \langle A;y.B(y)\rangle,\quad y\notin\operatorname{FV}(A).
\]
Bound variables are renamed before capture-free substitution.
The two variables $x,y$ in a paired formula are distinct.
An ordinary assumption $F$ is the constant clause
$\langle H;v.F\rangle$, where $v$ is absent from $F$ and $H$ is the
fixed computation theorem $I(0,0;u.0)=0$. Thus ordinary hypothesis
use is a clause Use preceded by this computation proof.

There are computation axioms and exactly the following inference
patterns:
\[
\frac{\Gamma\vdash A}{\Gamma\vdash B(t)}
\quad\mathsf{Use}
\quad(\langle A;y.B(y)\rangle\in\Gamma),
\]
\[
\frac{\Gamma\vdash\mathsf D_{x,y}B(x,y)\qquad
      \Gamma,\langle B(a,0);y.B(a,y)\rangle\vdash F}
     {\Gamma\vdash F}
\quad\mathsf{DE},
\qquad
\frac{\Gamma\vdash s=t\qquad\Gamma\vdash F(s)}
     {\Gamma\vdash F(t)}
\quad\mathsf{Tr},
\]
\[
\frac{\Gamma,B(t,\eta_\varepsilon(z))\vdash B(t,Sz)}
     {\Gamma\vdash\mathsf D_{x,y}B(x,y)}
\quad\mathsf{DI}_{\varepsilon},
\qquad
\eta_0(z)=0,\quad\eta_1(z)=z.
\tag{Choice DI}\label{eq:choice-di}
\]
Here $\varepsilon$ is a certificate choice in $\{0,1\}$, and
$\varepsilon=1$ is permitted only when $B$ has a literal equation
root. An equation beneath a binder is not a literal equation.
Call $\varepsilon=0$ the \emph{case choice} and $\varepsilon=1$
the \emph{progress choice}. Both retain exactly one ordinary
formula assumption. Compound matrices have only the case choice.

In DI the eigenparameter $z$ is fresh for $\Gamma$, the witness
$t$, and every fixed parameter of $B$. In DE the eigenparameter
$a$ is fresh for $\Gamma$, the conclusion $F$, and every fixed
parameter of $B$. In transport, $F(v)$ is any fixed native
formula template. Displayed local assumptions and clauses are
discharged. There are no primitive reflexivity, coverage,
separation, ordinary logical, or additional induction rules.

We call this calculus $\mathcal C$. The original comparison
calculus $\mathcal N$ has identical syntax, computation,
Use, DE, and transport, but its introduction is
\[
\frac{\Gamma,B(t,0),[B(t,z)]_{\mathrm{at}}\vdash B(t,Sz)}
     {\Gamma\vdash\mathsf D_{x,y}B(x,y)}.
\tag{Old DI}\label{eq:choice-old-di}
\]
The bracket supplies $B(t,z)$ exactly at a literal equation.

\subsection{Exact equivalence and the unused-witness construction}

\begin{lemma}[Structural admissibility]
\label{lem:choice-structural}
Weakening, exchange, contraction, renaming, ordinary-assumption cut,
and capture-free numerical substitution are admissible in $\mathcal C$.
\end{lemma}
\begin{proof}
Induct on derivations. Rename local eigenparameters before inserting
proofs or terms. Numerical substitution preserves a formula's root
constructor, so an allowed progress choice remains allowed.
The selected choice does not change under substitution or cut.
For ordinary-assumption cut, replace each use of its constant clause
by the supplied proof. Computation schemas are closed under
capture-free numerical substitution. These facts verify every case.
\end{proof}

\begin{theorem}[Identity-language equivalence]
\label{thm:choice-equivalence}
For every retained clause context $\Gamma$ and native formula $F$,
\[
 \Gamma\vdash_{\mathcal C}F
 \quad\Longleftrightarrow\quad
 \Gamma\vdash_{\mathcal N}F.
\]
The transformations are effective on finite derivations and leave
all numerical and formula syntax unchanged.
\end{theorem}
\begin{proof}
Each new introduction is an instance of the old one after weakening
by the omitted local assumption. For the case choice at an equation,
add the induction instance. For the progress choice, add the zero
instance. Compound case introduction is already the old rule.
Thus every new derivation translates to an old one.

For the reverse direction, transform an old derivation recursively.
Only its atomic DI case needs a construction. Suppose its matrix is
the equation $E(x,y)$ and its transformed premise is
\[
 \pi:\quad\Gamma,E(t,0),E(t,z)\vdash E(t,Sz).
\]
The outer new inference will be case DI with witness $t$. It is
therefore enough to prove $E(t,Sz)$ under
$\Gamma,E(t,0)$.

Rename the eigenvariable in $\pi$ to a fresh $v$. Choose a new
variable $u$ absent from $E(t,y)$ and all retained data. Under
$\Gamma,E(t,0)$, progress DI with witness zero introduces
\[
             \mathsf D_{u,y}E(t,y)
\]
from exactly the renamed premise $\pi[v/z]$.
The bound coordinate $u$ is deliberately unused.

Apply DE to this paired formula. Because $u$ is unused, the opened
clause is exactly
\[
               \langle E(t,0);y.E(t,y)\rangle;
\]
it contains no new existential witness. The retained ordinary
assumption $E(t,0)$ supplies its antecedent. Use at $Sz$ gives
$E(t,Sz)$, and DE closes. Finally the outer case DI introduces
the original $\mathsf D_{x,y}E(x,y)$ with witness $t$.
The original premise proof $\pi$ has been inserted once.

An old compound introduction already is the new case choice.
All other rule cases are unchanged. This completes the induction
on the original derivation.
\end{proof}

The unused inner coordinate is essential. Introducing
$\mathsf D_{x,y}E(x,y)$ internally and then eliminating it would
give an arbitrary fresh witness, not necessarily the original
$t$. The construction instead keeps $t$ as a fixed parameter
of $E(t,y)$ and binds only the unused coordinate $u$.

\begin{proposition}[Linear bound on inference nodes]
\label{prop:choice-node-bound}
Count each computation axiom, Use, DI, DE, and transport as one
inference node, expanding ordinary-assumption use as a computation
of $H$ followed by Use. If an old derivation has $N$ nodes and
$r$ atomic DI nodes, the displayed transformation has exactly
\[
                   N+5r\leq 6N
\]
nodes.
\end{proposition}
\begin{proof}
Each atomic old DI node is replaced by six nodes: outer case DI,
inner progress DI, DE, Use of the opened clause, Use of the
retained base assumption, and the computation of $H$ needed for
that last ordinary-assumption use. Its old premise subtree
occurs once. All other inference nodes are copied once.
\end{proof}

This counts inference-tree nodes. It is not a bound on serialized
JSON bytes, formula sizes, checker runtime, or the number of
consumed packet cells when hypothesis closure is implicit.
In the formula-context Q/D presentation, the corresponding
replacement has outer case DI, inner progress DI, DE, QE,
and two Hyp nodes, giving the same six-node count.

\subsection{Induction and numerical cases separate cleanly}

\begin{lemma}[Contextual equation induction]
\label{lem:choice-ei}
The progress choice, DE, and Use derive
\[
\frac{\Gamma\vdash E(0)\qquad
      \Gamma,E(z)\vdash E(Sz)}
     {\Gamma\vdash E(t)}
\tag{EI}\label{eq:choice-ei}
\]
for every equation $E$, with the usual freshness condition.
\end{lemma}
\begin{proof}
Choose an unused witness coordinate $u$. Progress DI introduces
$\mathsf D_{u,y}E(y)$ from the given step. DE opens the clause
$\langle E(0);y.E(y)\rangle$. Supply its antecedent using the given
base proof and use the clause at $t$. Close DE.
\end{proof}

\begin{lemma}[Cases, equality, and binder congruence]
\label{lem:choice-preinduction}
The case choice derives, for arbitrary native formulas,
\[
\frac{\Gamma\vdash F(0)\qquad\Gamma\vdash F(Sz)}
     {\Gamma\vdash F(t)}
\tag{NC}\label{eq:choice-nc}
\]
with the usual freshness condition. Equality reflexivity,
symmetry, transitivity, numerical congruence, and uniform
matrix interderivability under $\mathsf D$ are derivable
without using the progress choice.
\end{lemma}
\begin{proof}
For cases, use case DI with an unused witness coordinate and
matrix $F(y)$, ignoring its retained base in the successor
premise. DE opens $\langle F(0);y.F(y)\rangle$; the given zero
proof and Use at $t$ finish.

For reflexivity put $r=I(0,t;u.0)$. Use the computation equation
$r=t$ as both transport premises with motive $v=t$ to obtain
$t=t$. The usual explicit equality motives give symmetry,
transitivity, and numerical congruence.

For binder congruence, open the old paired formula at a fresh
witness $a$ and introduce the new one by case DI at $a$.
Convert its retained base to the old base, use the opened
old clause at the successor, and convert back. Ordinary cut
closes the construction, and the reverse conversion gives
the converse. No induction is used.
\end{proof}

The primitive progress choice is still restricted to equations.
Allowing it at arbitrary compound matrices would, by the
argument for \eqref{eq:choice-ei}, supply induction for
arbitrary formulas and change the intended theory.

\subsection{The entire ordinary logical interface}
\label{subsec:choice-logic}

Here is a predecessor-free reconstruction of the ordinary logical
operations. It makes explicit that the case choice supplies the
logical interface independently of progress induction.

Let $C(v,a,b):=I(v,a;u.b)$ with $u$ unused. Computation gives
$C(0,a,b)=a$ and $C(Sv,a,b)=b$. Define the gate on fixed finite
formula syntax by
\[
 G_v(s=t):=(C(v,0,s)=C(v,0,t)),\qquad
 G_v(\mathsf D_{x,y}B):=\mathsf D_{x,y}G_vB,
\]
renaming bound variables away from the control. Equality,
binder congruence, and external structural induction give
\[
 \vdash G_0A,\qquad G_{Sv}A\dashv\vdash A,\qquad
 A\dashv\vdash\mathsf D_{u,y}G_yA
 \quad(u,y\notin\operatorname{FV}(A)).
\tag{Gate and padding}\label{eq:choice-gate}
\]
For padding, case DI proves the forward direction. DE, the
zero-gate theorem, and Use at $1$ prove the reverse direction.

Define a selector between two arbitrary formulas by
\[
 M_v(s=t,u=w):=(C(v,s,u)=C(v,t,w))
\]
and, with aligned fresh paired variables, by
\[
\begin{aligned}
 M_v(\mathsf D_{x,y}A,\mathsf D_{x,y}B)
   &:=\mathsf D_{x,y}M_v(A,B),\\
 M_v(E,\mathsf D_{x,y}B)
   &:=\mathsf D_{x,y}M_v(G_yE,B),\\
 M_v(\mathsf D_{x,y}A,E)
   &:=\mathsf D_{x,y}M_v(A,G_yE),
\end{aligned}
\]
where $E$ is an equation. Every call decreases maximum input
binder depth. Padding and congruence prove
$M_0(A,B)\dashv\vdash A$ and
$M_{Sv}(A,B)\dashv\vdash B$.

Put $\delta(v)=C(v,1,0)$ and $p(z)=I(z,0;u.\delta(u))$.
Only the fixed computations $p(1)=1$ and $p(2)=0$ are needed.
The full logical macro inventory is
\[
\begin{aligned}
 A\Rightarrow_*B&:=\mathsf D_{u,z}M_z(A,B),\\
 A\wedge_*B&:=\mathsf D_{u,z}G_zM_{p(z)}(A,B),\\
 \exists_*x\,A(x)&:=\mathsf D_{x,y}G_yA(x),\\
 A\vee_*B&:=\exists_*x\,M_x(A,B),\\
 \forall_*x\,A(x)&:=
        A(0)\wedge_*\mathsf D_{u,z}G_zA(z),\\
 \bot_*&:=(0=1).
\end{aligned}
\tag{Logical macros}\label{eq:choice-macros}
\]
Newly bound variables are fresh; $u$ is unused.

Implication introduction obtains its assumption from the zero
selector and supplies the consequent at successors, then uses
case DI. Elimination opens its clause and uses it at $1$.
Existential introduction uses its witness and the successor
gate; elimination opens at a fresh witness, supplies the
zero-gate theorem, and uses the clause at $1$.

Disjunction introduction uses witnesses $0$ and $1$.
For elimination, discharge the two branch proofs as
$A\Rightarrow_*F$ and $B\Rightarrow_*F$. NC applied to
$M_z(A,B)\Rightarrow_*F$ gives this implication at every
control: the two endpoint cases follow from the branch
implications. Existential elimination provides $M_a(A,B)$
and the implication gives $F$.

Given $A$ and $B$, selector endpoints and NC give $M_v(A,B)$
at an arbitrary control. Substitute $v=p(Sz)$ and use the
successor gate for conjunction introduction. For projections,
open its zero-gated clause and specialize at $1$ and $2$;
the two fixed computations give $B$ and $A$, respectively.

These rules prove existential numerical coverage:
\[
     (x=0)\vee_*\exists_*v\,(x=Sv).
\tag{Coverage}\label{eq:choice-coverage}
\]
Use NC. Zero takes the left injection of reflexivity, and a
successor takes the right injection with the evident witness.
Coverage gives equality cases even when the retained context
contains $x$; its elimination adds $x=0$ or $x=Sv$ rather
than silently changing that retained context.

Universal introduction follows from a fresh-parameter proof
by substituting zero and successors and forming the two
conjuncts. For elimination, take its $A(0)$ conjunct and
open the paired conjunct. Its base is a zero-gate theorem.
Coverage of the target chooses the zero instance from
$A(0)$ or a successor instance from the opened clause
and gate removal.

Finally, under $0=1$, computation and transport give
$s=C(0,s,t)=C(1,s,t)=t$ for arbitrary numerical terms.
Every equation follows, and external induction on formula
syntax with case DI gives every native formula. This
supplies falsity elimination. No excluded middle for
quantified formulas has entered the construction.

\subsection{Whole-domain realizations and both formula composites}

For a standard first-order reading, let $T_{\rm eq}$ have
intuitionistic equality logic, the same numerical signature
and computation equations, numerical coverage, separation
$(0=1)\to\bot$, and contextual equation induction.
Its interpretation of the native syntax is
\[
 (s=t)^\circ=(s=t),\qquad
 (\mathsf D_{x,y}B)^\circ
   =\exists x\bigl(B^\circ(x,0)\to\forall y B^\circ(x,y)\bigr),
\]
\[
 \langle A;y.B\rangle^\circ
       =A^\circ\to\forall y B^\circ.
\]
The domain is the entire numerical domain and interpreted
equality is actual equality.

Case DI is sound in this realization by numerical coverage.
Progress DI is sound by contextual equation induction under
the conditional base. Use, DE, transport, and computation
have their displayed ordinary meanings. Conversely, the
native macros derive every ordinary logical rule,
\eqref{eq:choice-ei} derives equation induction, and
\eqref{eq:choice-coverage} derives coverage. Separation
translates to the identity implication
$(0=1)\Rightarrow_*(0=1)$.

For the essential paired formula composite,
\[
 \mathsf D_{x,y}B
 \dashv\vdash
 \exists_*x\bigl(B(x,0)\Rightarrow_*\forall_*y\,B(x,y)\bigr).
\tag{Paired composite}\label{eq:choice-paired-composite}
\]
Forwards, DE opens its clause at a fresh $a$; under $B(a,0)$
that clause supplies $B(a,y)$ at a fresh arbitrary $y$.
Introduce the universal, implication, and existential
macros and close DE. Backwards, eliminate the existential
at a fresh $a$. Introduce the paired formula with witness
$a$ by case DI: the retained base, supplied implication,
and universal elimination at $Sz$ give its successor premise.

The ordinary interpretations of the six logical macros
have their ordinary round trips. For conjunction, its true
base and the instances $1,2$ recover both operands;
the converse uses selector cases. The universal has all
successor instances in its paired conjunct and zero in
its first conjunct; coverage gives exactly the ordinary
universal. Implication, disjunction, and existential
quantification use their selector or gate endpoints.
Falsity uses the explicit ordinary separation bridge.
External induction on finite formulas, together with
\eqref{eq:choice-paired-composite}, gives both formula
composites. Clauses correspond in both directions to
$A\Rightarrow_*\forall_*yB$ by the derived logical rules.

The numerical composites are literal identities with
internally provable graph $x=y$. No standard-numeral
restriction, retract, quotient, or equality reflection
assumption is being inserted at this stage. The separately
supplied arithmetic bridge identifies $T_{\rm eq}$ with
full iPRA; the present change preserves that bridge
unchanged because Theorem~\ref{thm:choice-equivalence}
preserves every native sequent.

\subsection{What the theorem does and does not claim}

The new inference interface retains one ordinary assumption
at every paired introduction. It preserves the entire
consequence relation, not just equations or closed
theorems. Consequently the established full-iPRA
interpretations and equiconsistency transfer unchanged.

This result is different from retaining only progress
introduction and strengthening closure elimination.
Here the original elimination rule is retained, and
the case and progress introductions coexist. It is also
different from allowing arbitrary compound progress,
which would supply stronger induction.

A complete-checker comparison must still count the
certificate choice, its equation-root restriction,
substitution, context handling, computation, and the
public input/output interface. The mathematical theorem
alone is not a byte-count or minimality theorem.
The inference-node bound is distinct from source-size
measurements and from serialized certificate size.
The proof is a written universal proof; finite regression
tests do not constitute a Lean/Rocq formalization of it.

\section{Arithmetic strength and both ordinary interpretation composites}
\label{sec:ordinary-bridge}

We spell out the arithmetic bridge to avoid identifying standard numerical
agreement with bi-interpretability. First elaborate every fixed iterator
into an ordinary function symbol. At stage zero the symbols are $0,S$.
Given a fixed earlier-stage term $h(u,\vec p)$, introduce
$i_h(n,a,\vec p)$ with the two equations
\[
i_h(0,a,\vec p)=a,\qquad
i_h(Sn,a,\vec p)=h(i_h(n,a,\vec p),\vec p).
\]
Only finite descriptions occur; the domain contains no code for a
universal evaluator. Formal parameters are distinct when a description
is declared, while actual arguments can coincide.

Let $T_I$ be ordinary intuitionistic first-order logic with actual
equality and genuine falsity, these staged equations, contextual
induction for literal equations, and
\[
\forall x\,(x=0\vee\exists y\,x=Sy),
\qquad (0=S0)\to\bot.
\]
The second sentence is the falsity bridge. It must be stated: positive
numerical equations by themselves admit a singleton structure. Native
explosion from $0=S0$ and a translation of ordinary falsity to that
equation do not authorize silently dropping the bridge from $T_I$.

The native-to-ordinary translation is
\[
(s=t)^\circ=(s=t),\qquad
(\mathsf D_{x,y}B)^\circ=
\exists x\bigl(B^\circ(x,0)\to\forall y B^\circ(x,y)\bigr).
\]
A clause $\langle A;y.B\rangle$ becomes $A^\circ\to\forall yB^\circ$.
For the auxiliary closure used by the compact checkers,
$(\mathsf Q_yB)^\circ=B^\circ(0)\to\forall yB^\circ$.
The interpretation domain is all elements, defined by $x=x$;
equality is actual equality.

Case DI is sound because numerical coverage combines its retained base
with its successor proof. Progress DI is sound by literal-equation
induction under a temporarily assumed base. Use, DE and transport are
ordinary logical inferences under the translations. Conversely the
derived logical operations, equation induction, cases and explosion in
Section~\ref{sec:choice-introduction} recover the rules and axioms of
$T_I$. An auxiliary closure is interderivable with
$\mathsf D_{v,y}B(y)$ for unused $v$: in one direction use case DI with
the closure and its base; in the other eliminate the paired formula
and return its supplied closure. Thus formulas and context clauses have
the stated meanings in both presentations.

\subsection{The arithmetic dependency order}

The scalar update has no recursion-counter argument. We do not conceal
one in an encoding. The inherited internal proofs, reproduced in the
appendices, establish the required arithmetic in this order:
\begin{enumerate}
\item Using equation induction and numerical cases, construct a fixed
scalar predecessor term and prove $P0=0$ and $P(Sn)=n$.
Existential coverage then gives $x=0\vee x=S(Px)$: in the successor
case $x=Sy$, apply $P$ to obtain $Px=y$ and substitute.
\item Construct truncated subtraction and a numerical characteristic
of actual equality. Prove its reflection on the entire original domain.
\item Compile each fixed quantifier-free formula into its characteristic
equation and recover every contextual quantifier-free induction instance.
\item With that induction available, construct pairing and projections
by a bounded square-root search and compile every fixed general
primitive-recursive description into a scalar iterator term.
\end{enumerate}
No pairing or faithful equality detector is assumed to obtain
predecessor. No appeal to standard-model correctness replaces an
internal equality. Every induction fixes its passive parameters;
diagonal substitution occurs only into a previously proved identity.

For a target recursion
\[
f(0,\vec p)=g(\vec p),\qquad
f(Sn,\vec p)=h(n,f(n,\vec p),\vec p),
\]
let $\pi,L,H$ be the proved pairing and projections. The compiler uses
\[
V(n,\vec p)=I\bigl(n,\pi(0,g^*(\vec p));
u.\pi(SL(u),h^*(L(u),H(u),\vec p))\bigr),
\qquad f^*(n,\vec p)=H(V(n,\vec p)).
\]
Equation induction proves $L(V(n,\vec p))=n$; the two projection
identities then prove the target defining equations. The construction
is external recursion on each finite PR description, not an internal
quantification over arbitrary functions.

\begin{theorem}[Whole-domain bi-interpretability and equiconsistency]
The ordinary realization $T_I$ and full first-order iPRA are
bi-interpretable. In both theories $x=y$ defines the isomorphism from
the original domain to its double interpretation. The native Choice-DI
calculus is therefore a faithful proof presentation of a theory
bi-interpretable and equiconsistent with iPRA.
\end{theorem}
\begin{proof}
Interpret an iterator in iPRA by primitive recursion that ignores its
counter argument in the update. iPRA proves its equations, literal
equation induction, coverage and separation. In the reverse direction
use the scalar compiler. The predecessor and equality-reflection proofs
supply the successor axioms and every quantifier-free induction instance.
All ordinary logical rules are preserved.

For two iterators with the same seed and provably equal updates,
equation induction on the equality of their values proves uniqueness at
every counter, with all parameters fixed. The corresponding argument
applies to ordinary primitive recursion. Induct externally on a finite
description: its double translation satisfies its original defining
equations, so internal uniqueness identifies the two functions. This
also proves coherence when substitution changes the staged syntactic
description of an update.

The domain and equality are unchanged; hence the identity graph is
internally total, bijective and preserves every interpreted operation.
The displayed ordinary meanings of the logical macros and the paired
binder give both formula composites by external induction on finite
formula syntax. Clauses are similarly equivalent to their conditional
universal meanings. The explicit falsity bridge preserves contradiction
in both directions. Effective translations of finite proofs give
equiconsistency, without a claim that either theory proves its own
consistency.
\end{proof}

% Adapted with attribution from the pinned unparenthesized one-hole archive.
\section{Internal scalar predecessor bootstrap}
\label{sec:hybrid-bootstrap}

This scalar bootstrap is inherited from the supplied scalar development~\cite{erasure} and is reproduced to keep the scalar checker self-contained. It is not a new contribution of this paper. The proof uses a factor map fixing odd states and a saturated-iteration construction of predecessor.

The available arithmetic terms are $0$, $S$, and
\[
 I(n,a;u.h),\qquad
 I(0,a;u.h)=a,\qquad
 I(Sn,a;u.h)=h[I(n,a;u.h)/u].
\]
The update binds only $u$.  All other variables in it are fixed numerical
parameters.  We use equality reasoning, the already reconstructed
zero/successor cases, and induction on a single literal equation.
Every induction below keeps its other parameters fixed.  Predecessor,
successor injectivity, pairing, and induction on arbitrary
quantifier-free formulas are not assumptions of the argument.

\subsection{Elementary terms and two finite cycles}

Write $1=S0$ and use larger numerals as successor abbreviations.  Define
\[
 C(n,a,b)=I(n,a;u.b),\qquad
 Z(n)=C(n,1,0),\qquad
 a+b=I(b,a;u.Su).
\]
The update variable of $C$ is absent from its body.  Consequently
\[
 C(0,a,b)=a,\qquad C(Sn,a,b)=b.
\]
Numerical cases show that $Z(n)$ is $0$ or $1$.
The usual commutative-monoid equations for addition follow by literal
equation induction.  For example, $0+b=b$ and $Sa+b=S(a+b)$ are
proved by induction on $b$; these yield commutativity by a further
equation induction, and associativity follows by induction on the last
summand.  We use these equations to reorder finite sums.  The notation
$2a$ means $a+a$ and $3a$ means $a+a+a$; no multiplication operation
is needed for this bootstrap.

Put
\[
 N(n)=I(n,0;u.Z(u)),\qquad
 j(u)=Z(u)+2N(u),\qquad
 M(n)=I(n,0;u.j(u)).
\]
The range of $N$ is contained in $\{0,1\}$: its zero value is $0$,
and its successor value is a value of $Z$.  This uses numerical cases
on its counter, not induction on a disjunction.  Hence
\begin{equation}
 N(n)+Z(N(n))=1,\qquad
 Z(Z(N(n)))=N(n),\qquad N(n+2)=N(n).
 \label{eq:hybrid-parity}
\end{equation}
Here the last equation follows from two successor computations and
the preceding equation.  Literal-equation inductions give
\begin{equation}
 N(2n)=0,\qquad N(2n+1)=1,\qquad N(3n)=N(n).
 \label{eq:hybrid-parity-scaling}
\end{equation}
For the first equation the successor step adds two to the argument of
$N$; the second follows by one computation.  In the third, increasing
$n$ adds three to the left argument and therefore applies $Z$ once to
both sides of the induction equation.

Finite computations now give
\[
 j(0)=1,\qquad j(1)=2,\qquad j(2)=0.
\]
The proof below needs just these three values and
$M(Sn)=j(M(n))$; it does not assume a division algorithm or a theorem
characterizing $M$ as remainder modulo three.

\subsection{A ramp with a global scaling identity}

Define
\[
 G(n)=I(n,0;u.u+M(u)+1).
\]
The central invariant is
\begin{equation}
 M(G(n))=N(n).
 \label{eq:hybrid-ramp-invariant}
\end{equation}
Its zero case is immediate.  For its successor step, retain the
single induction equation and split the Boolean value $N(n)$.
If it is $0$, the next $G$ state is $G(n)+1$; its $M$ value is
$j(0)=1$, equal to $N(Sn)$.  If it is $1$, the next $G$ state
is $G(n)+2$; two computations of $M$ change its value from $1$
to $2$ and then to $0$, again equal to $N(Sn)$.  This proves
the invariant by literal-equation induction.

It follows that
\[
 G(Sn)=G(n)+N(n)+1,\qquad G(n+2)=G(n)+3.
\]
Equation induction therefore yields
\begin{equation}
 G(2n)=3n.
 \label{eq:hybrid-ramp-double}
\end{equation}
A separate literal-equation induction proves
\begin{equation}
 G(3n)=3G(n)+N(n).
 \label{eq:hybrid-ramp-triple}
\end{equation}
For completeness, its successor calculation is
\begin{align*}
 G(3Sn)
 &=G(3n)+N(3n)+4
   =3G(n)+2N(n)+4,\\
 3G(Sn)+N(Sn)
 &=3G(n)+3N(n)+3+Z(N(n)).
\end{align*}
The expressions agree in both Boolean cases for $N(n)$.
The base is $0=0$.  In particular,
\eqref{eq:hybrid-ramp-triple} was proved internally; it was not
inferred from an intended floor-function interpretation of $G$.

Set
\[
 H(u)=C(N(u),G(u),u).
\]
The crucial choice is to retain $u$ when its parity is $1$.
The preceding equations give
\begin{equation}
 H(2u)=3u,\qquad H(1)=1,\qquad H(3u)=3H(u).
 \label{eq:hybrid-factor-map}
\end{equation}
To check the last identity, use $N(3u)=N(u)$ and split the Boolean
$N(u)$.  At $0$, both sides are $3G(u)$ by
\eqref{eq:hybrid-ramp-triple}; at $1$, both are $3u$.
This argument does not require that an arbitrary number has already
been represented as $2v$ or $2v+1$.

\subsection{A comparator from scalar iteration}

For every fixed update $h$, equation induction on $n$, keeping $m,a$
and all parameters fixed, proves the semigroup identity
\[
 I(m+n,a;h)=I(n,I(m,a;h);h).
\]
The base is the zero computation and the step applies $h$ to the
induction equation. This identity will also be used in the later
pairing construction.


For a fixed update $h$, equation induction on $y$, with its initializer
$a$ and all parameters of $h$ fixed, proves the rotation identity
\begin{equation}
 I(Sy,a;h)=I(y,h(a);h).
 \label{eq:hybrid-rotation}
\end{equation}
Both sides have base $h(a)$, and their successor equations apply the
same $h$ to the previous equality.  Here and below $I(y,a;h)$ is
an abbreviation for one fixed finite scalar term, not quantification
over functions.

By \eqref{eq:hybrid-factor-map}, another equation induction gives
\begin{equation}
 I(y,3a;H)=3I(y,a;H).
 \label{eq:hybrid-iterator-scaling}
\end{equation}
Its base is $3a=3a$.  In the step the fixed-parameter induction
equation rewrites the new left side as
\[
 H(3I(y,a;H))=3H(I(y,a;H)).
\]
The initializer parameter stays $a$ throughout this induction.
Equation induction also gives $I(y,1;H)=1$, since $H(1)=1$.

Define
\[
 E(x)=I(x,1;u.2u),\qquad
 L(x,y)=N\bigl(I(y,E(x);H)\bigr).
\]
Then the following entire recursion specification is provable:
\begin{equation}
 L(0,y)=1,\qquad L(Sx,0)=0,\qquad L(Sx,Sy)=L(x,y).
 \label{eq:hybrid-comparator}
\end{equation}
The first equation uses $E(0)=1$ and the fixed-point identity for
$H$.  The second uses $E(Sx)=2E(x)$ and $N(2a)=0$.
For the third,
\begin{align*}
 L(Sx,Sy)
 &=N\bigl(I(Sy,2E(x);H)\bigr)\\
 &=N\bigl(I(y,H(2E(x));H)\bigr)\\
 &=N\bigl(I(y,3E(x);H)\bigr)\\
 &=N\bigl(3I(y,E(x);H)\bigr)\\
 &=N\bigl(I(y,E(x);H)\bigr)=L(x,y).
\end{align*}
The successive steps use rotation, $H(2a)=3a$, iterator scaling,
and $N(3a)=N(a)$.  Thus $L$ has the comparison equations at all
arguments internally.  Its standard-model meaning as the indicator
of $x\le y$ is motivation, not a premise of the following argument.

\subsection{Saturated iteration gives predecessor}

We now need only the three equations
\eqref{eq:hybrid-comparator}.  Define
\[
 B_m(u)=C(L(Su,m),u,Su),\qquad
 W(t,m)=I(t,0;u.B_m(u)).
\]
The subscript in $B_m$ records a fixed parameter; $B_m$ is an open
term abbreviation, not a function variable.

Numerical cases on $e$ prove the selector equation
\[
 C(e,Sa,Sb)=S C(e,a,b).
\]
Equivalently one may use equation induction on $e$, whose induction
hypothesis is unused.  Therefore
\begin{align}
 B_{Sm}(Su)
 &=C(L(SSu,Sm),Su,SSu)\notag\\
 &=C(L(Su,m),Su,SSu)=S B_m(u),
 \label{eq:hybrid-conjugacy}\\
 B_{Sm}(0)&=C(L(S0,Sm),0,S0)=C(L(0,m),0,S0)=S0.
 \label{eq:hybrid-conjugacy-base}
\end{align}

Literal-equation induction on $t$, with $m$ fixed, proves
\begin{equation}
 W(St,Sm)=S W(t,m).
 \label{eq:hybrid-saturated-shift}
\end{equation}
The base is \eqref{eq:hybrid-conjugacy-base}.  The successor step is
\begin{align*}
 W(SSt,Sm)
 &=B_{Sm}(W(St,Sm))\\
 &=B_{Sm}(S W(t,m))\\
 &=S B_m(W(t,m))=S W(St,m).
\end{align*}
Only the second line uses the current induction hypothesis, at the
same fixed $m$.

Two subsequent equation inductions give
\begin{equation}
 L(Sn,n)=0,\qquad W(n,n)=n.
 \label{eq:hybrid-diagonals}
\end{equation}
The first uses the comparator recurrence.  The second uses
\eqref{eq:hybrid-saturated-shift} in its successor step.  The
substitution $t=m=n$ is made into that already proved universally
parameterized equation; no parameter of an active induction hypothesis
is changed.  Consequently
\begin{equation}
 W(Sn,n)=B_n(n)=C(L(Sn,n),n,Sn)=n.
 \label{eq:hybrid-extra-step}
\end{equation}

Finally define the finite scalar term
\begin{equation}
 P(m)=I\bigl(m,0;u.C(L(SSu,m),u,Su)\bigr).
 \label{eq:hybrid-predecessor-term}
\end{equation}
Its zero equation is immediate.  To justify its successor equation
without assuming a separate extensionality rule for iterator bodies,
put
\[
 V(t,m)=I\bigl(t,0;u.C(L(SSu,Sm),u,Su)\bigr).
\]
The comparator recurrence identifies this update pointwise with
$B_m(u)$.  Equation induction on $t$, with $m$ fixed, gives
$V(t,m)=W(t,m)$: their base values are $0$, and the successor step
uses the pointwise update equation and then the induction equation.
Now \eqref{eq:hybrid-extra-step} yields
\[
 \boxed{\qquad P(0)=0,\qquad P(Sm)=V(Sm,m)=W(Sm,m)=m.\qquad}
\]

\subsection{Dependency boundary}

All terms in this section are built from $0,S,I$ by finite substitution.
All inductions have a single equation as their invariant.  Numerical
cases and equality transport are the already reconstructed logical
resources.  Successor injectivity may now be obtained by applying $P$
to an assumed equation $Sx=Sy$.  Equality tests, reflection of
quantifier-free formulas, pairing, and the recovery of general
primitive recursion can then be developed in the established order.
They were not used in order to manufacture predecessor.

This argument is a written internal derivation in the stated theory.
The finite implementation checks do not constitute a proof-assistant
formalization of this entire metatheorem.

% Standalone section.  Requires only amsmath and amssymb in the parent document.
% No theorem environments, macros, or bibliography keys are assumed.
% All arithmetic claims in this section are proved here.

\section{From literal-equation induction to full intuitionistic PRA}
\label{sec:eic-completion}\label{app:arithmetic}

This section supplies an unconditional completion theorem. It is adapted
with attribution from the current forward-left draft~\cite{forward}. In the
scalar application used here, every function occurring in the proof is
already a scalar term: predecessor is supplied by
Appendix~\ref{sec:hybrid-bootstrap}, the selector is
$C(n,a,b)=I(n,a;u.b)$, addition is $x+y=I(y,x;u.Su)$, and
truncated subtraction is $d(x,y)=I(y,x;u.Pu)$. All other functions below
are finite compositions of these. Thus this proof does not presuppose
that arbitrary primitive recursion has already been compiled. The QFI
conclusion first applies to the scalar language; the general PR compiler
comes afterward.  Literal-equation
induction, together with the primitive-recursive defining equations,
predecessor coverage, and separation of $0$ and $1$, proves induction for every
quantifier-free formula.  In particular, the proof does not assume a faithful
numerical equality test, stability or decidability of equality, successor
uniqueness, or any commutation law for an iterator.  We construct the equality
test and prove its reflection property for arbitrary elements of the
numerical domain.

\subsection{The exact hypotheses}

Work in ordinary intuitionistic first-order logic with equality and genuine
falsity $\bot$.  The language has one numerical sort and a symbol for each
fixed, finite primitive-recursive description.  Distinct descriptions are
not identified merely because they compute the same standard function.
Let $\mathsf{PRE}$ contain the defining equations for zero, successor,
projections, composition, and primitive recursion.  In particular, choose
descriptions $P$ and $C$ with equations
\begin{align}
 P(0)&=0, & P(Sx)&=x, \label{eq:eic-predecessor}\\
 C(0,a,b)&=a, & C(Sx,a,b)&=b. \label{eq:eic-selector}
\end{align}
Write $1=S0$.  In addition to $\mathsf{PRE}$, assume
\begin{align}
 \mathsf{Cov}:&\qquad
       \forall x\,\bigl(x=0\lor x=S(Px)\bigr),
       \label{eq:eic-coverage}\\
 \mathsf{Sep}:&\qquad 0=1\longrightarrow\bot,
       \label{eq:eic-separation}
\end{align}
and induction for each \emph{literal equation} $E(n,\vec p)$:
\begin{equation}
 \forall\vec p\,
 \left[
   \left(E(0,\vec p)\land
       \forall n\bigl(E(n,\vec p)\longrightarrow E(Sn,\vec p)\bigr)
   \right)
   \longrightarrow\forall n\,E(n,\vec p)
 \right].
 \label{eq:eic-induction}
\end{equation}
Here the schematic expression $E$ is a single equality of numerical
terms; the displayed universal quantifiers range only over numbers.
Every description and every instance of the schema is finite external
syntax.

The axiom formulation \eqref{eq:eic-induction} and the usual contextual
induction rule are interchangeable.  For example, if
\[
 \Gamma\vdash E(0,\vec p),
 \qquad
 \Gamma,E(z,\vec p)\vdash E(Sz,\vec p),
\]
where $z$ is absent from $\Gamma$ and the passive parameters, discharge the
second premise's equation and universally introduce $z$.  The axiom then
gives $\Gamma\vdash\forall n E(n,\vec p)$, which can be instantiated at any
target term.  Conversely, assume the antecedent of the axiom and apply the
contextual rule with fresh variables.  This uses the ordinary intuitionistic
logical rules and permits arbitrary, including quantified, formulas in
$\Gamma$.

\subsection{Constructor separation and two elementary tests}

Define the sign description and the countdown description by
\begin{equation}
 \operatorname{sg}(u)=C(u,0,1),
 \qquad
 K(u)=C(u,1,Pu).
 \label{eq:eic-sign-countdown}
\end{equation}
Their defining equations imply
\begin{align}
 \operatorname{sg}(0)&=0,
 &\operatorname{sg}(Sx)&=1,\label{eq:eic-sign-computation}\\
 K(0)&=1,
 &K(Sx)&=x.\label{eq:eic-countdown-computation}
\end{align}

\paragraph{Constructor laws.}
Successor is injective: applying $P$ to $Sx=Sy$ gives $x=y$ by
\eqref{eq:eic-predecessor}.  Successor never equals zero.  Indeed,
$Sx=0$ gives $1=0$ by applying $\operatorname{sg}$ and using
\eqref{eq:eic-sign-computation}; symmetry and
\eqref{eq:eic-separation} give $\bot$.  Consequently $1\ne0$, where
$A\ne B$ abbreviates $A=B\to\bot$.  None of these arguments uses
induction.

\paragraph{The sign test.}
For every $u$,
\begin{align}
 &\operatorname{sg}(u)=0\lor\operatorname{sg}(u)=1,
       \label{eq:eic-sign-range}\\
 &\operatorname{sg}(u)=0\quad\longleftrightarrow\quad u=0,
       \label{eq:eic-sign-zero}\\
 &\operatorname{sg}(u)=1\quad\longleftrightarrow\quad u\ne0.
       \label{eq:eic-sign-one}
\end{align}
To prove the range assertion, split \eqref{eq:eic-coverage}.  The branch
$u=0$ gives sign $0$, and the branch $u=S(Pu)$ gives sign $1$.
The implication $u=0\to\operatorname{sg}(u)=0$ is computation.
For its converse, the successor branch of coverage would make
$\operatorname{sg}(u)=0$ imply $1=0$, a contradiction; in that branch
explosion gives $u=0$ as well.  Similarly, if
$\operatorname{sg}(u)=1$ and $u=0$, computation gives $0=1$.
Conversely, $u\ne0$ excludes the zero branch of coverage, and the
successor branch computes the sign as $1$.

Coverage and \eqref{eq:eic-selector} also give, for all $u,v$,
\begin{equation}
 C(u,v,v)=v.
 \label{eq:eic-selector-idempotence}
\end{equation}
Both coverage branches compute to $v$, so this identity again needs no
induction.

\paragraph{The countdown test and the notation $W$.}
For every $u$,
\begin{equation}
 K(u)=0\quad\longleftrightarrow\quad u=1.
 \label{eq:eic-countdown-zero}
\end{equation}
If $u=1=S0$, \eqref{eq:eic-countdown-computation} gives $K(u)=0$.
Conversely, assume $K(u)=0$ and split coverage.  If $u=0$, then
$K(u)=1$, contradicting separation, and explosion gives $u=1$.
If $u=S(Pu)$, computation gives $K(u)=Pu$; hence $Pu=0$ and
$u=S0=1$.

If a presentation uses a description named $W$ with defining equations
\[
 W(0)=1,\qquad W(Sx)=x,
\]
then it proves $K(u)=W(u)$ for every $u$ already from coverage.
In the zero branch both sides are $1$; in the branch $u=S(Pu)$ both
sides are $Pu$.  Thus $K$ and $W$ are provably equal on the whole domain,
even if they are distinct descriptions.  This equality does not appeal to
standard-model agreement or to uniqueness of an arbitrary recurrence.

\subsection{Four elementary uses of equation induction}

Choose descriptions for addition and truncated subtraction, with recursion
on the second argument:
\begin{align}
 x+0&=x, &x+Sy&=S(x+y),\label{eq:eic-addition}\\
 d(x,0)&=x, &d(x,Sy)&=P(d(x,y)).\label{eq:eic-subtraction}
\end{align}
Using these equations and \eqref{eq:eic-induction}, we now derive
\begin{align}
 0+y&=y,\label{eq:eic-left-zero}\\
 (Sx)+y&=S(x+y),\label{eq:eic-left-successor}\\
 d(Sx,Sy)&=d(x,y),\label{eq:eic-cancel-successors}\\
 d(Sa,a)&=1.\label{eq:eic-diagonal}
\end{align}
Each induction below has a single equation as its induction formula.

For \eqref{eq:eic-left-zero}, induct on $y$.  The base is $0+0=0$.
The step is $0+Sy=S(0+y)=Sy$.

For \eqref{eq:eic-left-successor}, hold $x$ fixed and induct on $y$.
At zero both sides are $Sx$.  At a successor the induction hypothesis
gives
\[
 (Sx)+Sy=S\bigl((Sx)+y\bigr)
       =S\bigl(S(x+y)\bigr)=S(x+Sy).
\]
Combining \eqref{eq:eic-left-successor} with the defining equation for
addition gives the useful rearrangement
\begin{equation}
 (Sy)+r=y+Sr.
 \label{eq:eic-shift-successor}
\end{equation}
Indeed, both sides are $S(y+r)$.

For \eqref{eq:eic-cancel-successors}, hold $x$ fixed and induct on $y$.
The base computes as
\[
 d(Sx,S0)=P(d(Sx,0))=P(Sx)=x=d(x,0).
\]
Applying $P$ to the induction equation gives the step:
\[
 d(Sx,S(Sy))=P(d(Sx,Sy))=P(d(x,y))=d(x,Sy).
\]

Finally, induct on $a$ for \eqref{eq:eic-diagonal}.  The base is
$d(S0,0)=S0=1$.  At a successor,
\[
 d(S(Sa),Sa)=d(Sa,a)=1,
\]
using \eqref{eq:eic-cancel-successors} and the induction hypothesis.

\subsection{An equation invariant that reflects actual equality}

The remaining arithmetic ingredient is the following equation, valid for
all $x,y$:
\begin{equation}
 C\bigl(\operatorname{sg}(d(x,y)),\ x,\ y+d(x,y)\bigr)=x.
 \label{eq:eic-residual-invariant}
\end{equation}
Its meaning is that a nonzero subtraction residual reconstructs the
original argument by addition.  The selector expresses this assertion as
one literal equation, so equation induction applies directly.

Hold $x$ fixed and induct on $y$.  At $y=0$, subtraction gives residual
$x$, and \eqref{eq:eic-left-zero} gives $0+x=x$.  The desired equation
is consequently $C(\operatorname{sg}(x),x,x)=x$, which is
\eqref{eq:eic-selector-idempotence}.

For the step, assume \eqref{eq:eic-residual-invariant} at $y$ and
abbreviate
\[
 r=d(x,y),\qquad r'=Pr=d(x,Sy).
\]
These are local term abbreviations, not additional function symbols.
Split coverage for $r'$.

If $r'=0$, then $\operatorname{sg}(r')=0$, and the successor instance
of \eqref{eq:eic-residual-invariant} reduces immediately to $x=x$.

In the other branch, $r'=S(Pr')$, so constructor separation gives
$r'\ne0$.  Split coverage for $r$.  Its zero branch would imply
$r'=P0=0$, a contradiction, from which the desired step equation follows.
In its successor branch,
\[
 r=S(Pr)=Sr',\qquad \operatorname{sg}(r)=1.
\]
The induction equation therefore computes to $y+r=x$.  Meanwhile,
\eqref{eq:eic-shift-successor} gives
\[
 (Sy)+r'=y+Sr'=y+r=x.
\]
Since $r'$ is a successor, $\operatorname{sg}(r')=1$, so this is
exactly the branch selected in the successor instance of
\eqref{eq:eic-residual-invariant}.  This completes the equation induction.
All case distinctions came from coverage; no induction on an implication
or a disjunction was used.

\paragraph{Reflection at residual $1$.}
We have, for arbitrary $a,b$,
\begin{equation}
 d(Sa,b)=1\quad\longleftrightarrow\quad a=b.
 \label{eq:eic-residual-reflection}
\end{equation}
For the forward direction, substitute $x=Sa$, $y=b$ into
\eqref{eq:eic-residual-invariant}.  If $d(Sa,b)=1$, the sign is $1$
and the selector returns its second branch, giving
\[
 b+1=Sa.
\]
The defining equation for addition gives $b+1=S(b+0)=Sb$.  Thus
$Sb=Sa$, and applying $P$ gives $b=a$, hence $a=b$.
Conversely, $a=b$ and \eqref{eq:eic-diagonal} give
$d(Sa,b)=d(Sa,a)=1$.

Define the equality characteristic by the fixed description
\begin{equation}
 e(a,b)=\operatorname{sg}\bigl(K(d(Sa,b))\bigr).
 \label{eq:eic-equality-characteristic}
\end{equation}
The sign, countdown, and residual reflection lemmas give
\begin{align}
 e(a,b)=0&\quad\longleftrightarrow\quad a=b,
       \label{eq:eic-equality-reflection}\\
 e(a,b)=0&\ \lor\ e(a,b)=1.
       \label{eq:eic-equality-range}
\end{align}
In detail, $e(a,b)=0$ is equivalent to $K(d(Sa,b))=0$ by
\eqref{eq:eic-sign-zero}, which is equivalent to $d(Sa,b)=1$ by
\eqref{eq:eic-countdown-zero}, and then to $a=b$ by
\eqref{eq:eic-residual-reflection}.  The range assertion is
\eqref{eq:eic-sign-range} applied to $K(d(Sa,b))$.

The reflection property concerns the logic's actual equality throughout
the domain.  It was not assumed as part of a numerical comparison symbol,
and its proof did not restrict $a$ or $b$ to standard numerals.
It also yields decidable equality.  In the $e(a,b)=0$ branch,
\eqref{eq:eic-equality-reflection} gives $a=b$.  In the $e(a,b)=1$
branch, an assumption $a=b$ would give $e(a,b)=0$ and therefore $0=1$,
contradicting separation.  Hence
\begin{equation}
 \forall a\,\forall b\,\bigl(a=b\lor(a=b\to\bot)\bigr)
 \label{eq:eic-equality-decidable}
\end{equation}
is a theorem, not an extra hypothesis.

\subsection{Compilation of all quantifier-free formulas}

For each fixed quantifier-free formula $A(\vec p)$, construct a term
$b_A(\vec p)$ by external recursion on its finite syntax.  Take equality,
$\bot$, $\land$, $\lor$, and $\to$ as the primitive formula constructors;
negation is implication into $\bot$.  Set
\begin{align}
 b_{s=t}&=e(s,t), & b_{\bot}&=1,\label{eq:eic-boolean-atoms}\\
 b_{A\land D}&=C(b_A,b_D,1),\label{eq:eic-boolean-and}\\
 b_{A\lor D}&=C(b_A,0,b_D),\label{eq:eic-boolean-or}\\
 b_{A\to D}&=C(b_A,b_D,0).\label{eq:eic-boolean-implies}
\end{align}
Every term here is a finite composition of fixed primitive-recursive
descriptions; if desired, projections and composition give it a single
description symbol.  No internal quantification over formulas or programs
is involved.

An external induction on the syntax of $A$ proves, within the arithmetic
theory,
\begin{equation}
 \bigl(b_A=0\lor b_A=1\bigr)
 \quad\text{and}\quad
 \bigl(b_A=0\longleftrightarrow A\bigr).
 \label{eq:eic-boolean-correctness}
\end{equation}
Here are the full connective checks.  In every case $b_A=1$ implies
$\neg A$: if $A$ held, the already established implication
$A\to b_A=0$ would contradict separation.

\paragraph{Atomic formulas and falsity.}
For an equation, both claims are
\eqref{eq:eic-equality-reflection}--\eqref{eq:eic-equality-range}.
For falsity the range assertion is $1=0\lor1=1$, with its right disjunct
proved by reflexivity.  Also $1=0\to\bot$ follows from separation and
symmetry, while $\bot\to1=0$ follows by explosion.  Thus $1=0$ is
equivalent to $\bot$, as required.

\paragraph{Conjunction.}
Split $b_A=0\lor b_A=1$.  In the zero branch, $A$ holds and
$b_{A\land D}=b_D$.  Its range follows from the range of $b_D$,
and its being zero is equivalent to $D$, hence to $A\land D$ under
the known $A$.  In the one branch, $\neg A$ holds and
$b_{A\land D}=1$.  Its range is immediate.  If it were zero,
separation would give $\bot$, hence $A\land D$; conversely an
assumption $A\land D$ contradicts $\neg A$, and explosion gives
$b_{A\land D}=0$.

\paragraph{Disjunction.}
In the branch $b_A=0$, the formula $A$ holds and
$b_{A\lor D}=0$.  The range is immediate, and both $b_{A\lor D}=0$
and $A\lor D$ hold, proving their equivalence.  In the branch
$b_A=1$, we have $\neg A$ and $b_{A\lor D}=b_D$.  Its range follows
from the range of $b_D$.  If its value is zero, then $D$ and hence
$A\lor D$.  Conversely, from $A\lor D$, the $A$ alternative
contradicts $\neg A$ and yields $b_D=0$ by explosion, while the
$D$ alternative yields $b_D=0$ directly.  This proves the required
equivalence in both alternatives.

\paragraph{Implication.}
In the branch $b_A=0$, the formula $A$ holds and
$b_{A\to D}=b_D$.  Its range follows from the range of $b_D$.
If its value is zero, then $D$, which gives $A\to D$.  Conversely,
$A\to D$ and the known $A$ give $D$, hence $b_D=0$.
In the branch $b_A=1$, we have $\neg A$ and
$b_{A\to D}=0$.  Its range is immediate.  To prove $A\to D$,
assume $A$; this contradicts $\neg A$, and explosion gives $D$.
Thus both $b_{A\to D}=0$ and $A\to D$ hold in this branch, and
their equivalence follows.  This argument uses the proved two-value
range of the finite characteristic.  It does not use excluded middle
for arbitrary formulas or change the surrounding intuitionistic logic.

This completes the external syntactic induction proving
\eqref{eq:eic-boolean-correctness} for every quantifier-free formula.

\subsection{Recovery of induction and the full theory}

Let $A(n,\vec p)$ be any fixed quantifier-free formula.  Suppose an
arbitrary retained context $\Gamma$ proves its base and its induction
step:
\[
 \Gamma\vdash A(0,\vec p),
 \qquad
 \Gamma,A(z,\vec p)\vdash A(Sz,\vec p),
\]
with the usual freshness condition on $z$.  By
\eqref{eq:eic-boolean-correctness}, the base gives
$\Gamma\vdash b_A(0,\vec p)=0$.  Under the additional assumption
$b_A(z,\vec p)=0$, the same equivalence gives $A(z,\vec p)$,
the given step gives $A(Sz,\vec p)$, and the equivalence gives
$b_A(Sz,\vec p)=0$.  Apply literal-equation induction to
$b_A(n,\vec p)=0$ and then apply its reflection implication.
The result is $\Gamma\vdash A(t,\vec p)$ for every target term $t$.
Thus the theory proves the full contextual quantifier-free induction
rule, including in arbitrary quantified contexts.

\paragraph{Completion theorem.}
In the common language of fixed primitive-recursive descriptions, the
theory
\[
 \mathsf{PRE}+\mathsf{Cov}+\mathsf{Sep}+\mathsf{EI}
\]
has exactly the same sequents as ordinary first-order intuitionistic
primitive recursive arithmetic with all quantifier-free induction
instances.  The forward simulation follows from the recovered
quantifier-free induction, successor injectivity, and successor separation;
all primitive-recursive defining equations and all logical rules are
shared.  For the reverse simulation, literal-equation induction is an
instance of quantifier-free induction, and successor separation gives
$\mathsf{Sep}$.  Quantifier-free induction proves $\mathsf{Cov}$:
at zero choose its left disjunct, and at a successor use
$Sx=S(P(Sx))$, which follows from $P(Sx)=x$, for its right disjunct.

Both interpretations therefore use the full domain, defined by $x=x$,
and preserve every function symbol, equality, logical connective, and
quantifier.  Both composites have the identity isomorphism, with graph
$x=y$.  The theories are consequently identity-domain bi-interpretable
and equiconsistent.  This completion theorem is unconditional once the
literal-equation induction rule has been obtained; it has no dependency
on a classical successor-uniqueness reduction.

\section{Pairing, primitive recursion, and internal numerical composites}
\label{sec:iterator-coding}
This appendix is adapted with attribution from the progress-closure
arithmetic development~\cite{progress}. Its prerequisites have now been
proved: scalar predecessor, arithmetic reflection, and full quantifier-free
induction. Pairing is a derived proof tool; it is not part of any checker.

\subsection{Pairing from a bounded square-root search}

Once QFI is available, a familiar square-based injection gives a
particularly direct compiler. Write $a\le b$ for $d(a,b)=0$.
The usual linear-order and semiring identities follow by fixed
quantifier-free inductions from the successor, addition, and monus
equations. Multiplication itself is the scalar term
\[
                        a\cdot b=I(b,0;u.u+a).
\]
Only primitive-recursive arithmetic is involved in the following
bounded root search; there is no minimization over an unbounded domain.

For a fixed parameter $n$, put
\[
 U_n(v)=C(d((Sv)^2,n),Sv,v),\qquad
 R_n(i)=I(i,0;v.U_n(v)),\qquad q(n)=R_n(n).
\]
The update increments $v$ exactly when $(v+1)^2\le n$, and
otherwise retains it. QFI in $i$, with $n$ fixed, proves
\[
 R_n(i)\le i,\qquad R_n(i)^2\le n,\qquad
 R_n(i)=i\ \vee\ n<(S(R_n(i)))^2.              \tag{Root invariant}
\]
The base has $R_n(0)=0$. If the update increments, its test
contradicts the strict-upper-bound alternative in the third invariant.
Thus the old state equals $i$, and the new state equals $Si$;
its square bound is the increment test. If the update does not increment, the
strict upper square bound holds and the state remains fixed.
An already established strict upper square bound forces every
subsequent update to remain fixed. These are exactly the cases
needed for the three displayed invariants.

At $i=n$, either the strict upper square bound already holds or
$R_n(n)=n$, in which case $n<(n+1)^2$ supplies it. Thus
\[
                        q(n)^2\le n<(q(n)+1)^2.           \tag{Root}
\]
Monotonicity of squares shows uniqueness of a number satisfying
Root. All formulas just used are quantifier free in the scalar
language, with fixed $n$ during the induction.

Now define
\[
 \pi(a,b)=(a+b)^2+a,\qquad
 L(n)=d(n,q(n)^2),\qquad H(n)=d(q(n),L(n)).
\]
For $k=a+b$, one has $k^2\le k^2+a<(k+1)^2$ since $a\le k$.
Root uniqueness gives $q(\pi(a,b))=a+b$, and therefore
\[
                    L(\pi(a,b))=a,\qquad H(\pi(a,b))=b.
\]
Surjectivity of $\pi$ is unnecessary. The original numerical domain
is never replaced by the range of this coding; pairing is used only
to store the finite state of a fixed primitive recursion. This
standard square pairing is a convenient reconstruction here, not
a claimed new arithmetic basis.

\subsection{Compiling each primitive recursion}

Suppose the fixed target description is
\[
 f(0,\vec p)=g(\vec p),\qquad
 f(Sn,\vec p)=h(n,f(n,\vec p),\vec p).
\]
Assume by external structural induction that the earlier $g,h$
already have scalar translations $g^*,h^*$. Define
\[
 W(u,\vec p)=\pi(SL(u),h^*(L(u),H(u),\vec p)),
\]
\[
 V(n,\vec p)=I(n,\pi(0,g^*(\vec p));u.W(u,\vec p)),
 \qquad f^*(n,\vec p)=H(V(n,\vec p)).
\]
The update has only its state variable bound, with every $p_i$
fixed. Equation induction proves $L(V(n,\vec p))=n$.
The two projection equations and scalar computation then prove
both defining equations for $f^*$ internally. Initial functions
and composition translate directly. This completes the translation
of every fixed finite PR description without adding a universal
evaluator.

\subsection{Why both numerical round trips are internal}

Two fixed scalar iterators with the same seed and provably equal
updates are internally equal at all counters: use literal-equation
induction comparing their values. Their base agrees, and their
successor equations plus congruence prove the step. The same
uniqueness argument works for ordinary primitive recursion.

For an ordinary PR description, induct externally on its finite
construction. The double translation satisfies its original
defining equations, so internal equation induction identifies it
with the original operation. For a native scalar description,
earlier translated updates agree by the external induction
hypothesis; scalar uniqueness identifies the double-translated
iterator. These proofs hold at arbitrary internal counters.

Thus the domain relation $x=y$ is provably a total bijection
preserving every fixed numerical operation in both composites.
The formula constructions and round trips in
Section~\ref{sec:choice-introduction} give the corresponding formula
composites. Both interpretations preserve falsity, taking native
$0=1$ to an ordinary contradiction and ordinary $\bot$ to native
$0=1$. Finite proof translations therefore establish equiconsistency.

\clearpage
\section{The recommended standalone checker}
\label{app:code}

This is the complete executable module, including both input checks and
the public verdict. The separately supplied readable expansion explains
the same algorithm with descriptive names; it is independently executable.
The code below uses the symmetric administrative triple reading from
Section~\ref{sec:raw-reversal}.
\begin{lstlisting}[language=Python]
import json

def c(a, g=[0], k=0):
    i = lambda a, t=c, d=0: [i(b, t, d + 5 - len(a)) for b in a] if a * 0 == [] else a if k <= a < d else d % 2 and i(t, d=3)
    if i(a) in g:
        return
    m, s, u, t = q()
    a = m == ([s, u] if s else [t] * 3) and i(t, u) == a[::-1] and i(t, s) or a
    h = c(m, g, k) or m * 3 or (u and [[s and [s], u, t], i(t, [s, u, t], 3) if s else u]) or (a := i(*a, t, 2))
    k -= 1
    c(i(a, [k]), g + [i(h, (m or len(a) % 2 == u) and k)], k)

def v(a, d=0):
    if type(a) is int:
        return d % 2 and (a == 0 or (3 <= a < d and a % 2))
    return type(a) is list and len(a) in ([1, 3] if d % 2 else [1, 2, 3]) and (d % 2 or len(a) != 3 or a == [a[0]] * 3) and all((v(b, d + 5 - len(a)) for b in a))

def j(a):
    return type(a) is int or (type(a) is list and all((j(b) for b in a)))
q = json.load(open(0))
if not j(q):
    raise ValueError('Only integer/list certificate syntax is allowed')
q = q.pop
a = q()
if not v(a):
    raise ValueError('The goal is not a canonical closed formula')
c(a)
print('ACCEPT')
\end{lstlisting}

\begin{thebibliography}{99}
\footnotesize
\bibitem{baseline} ProveIt contributors,
\emph{Concise presentations of intuitionistic primitive recursive arithmetic},
README at commit \code{1517d9d79ce33f631e0f82960c849427cbba23bf}.
\href{https://github.com/VladimirReshetnikov/ProveIt/blob/1517d9d79ce33f631e0f82960c849427cbba23bf/Logic/PrimitiveRecursiveArithmetic/ConciseIPRA/README.md}{Pinned repository README}.
\bibitem{leanbridge} ProveIt contributors,
\emph{NativeBiInterpretation.lean, NativeEq.lean, and BiInterpretation.lean},
same pinned commit.
\href{https://github.com/VladimirReshetnikov/ProveIt/tree/1517d9d79ce33f631e0f82960c849427cbba23bf/Logic/PrimitiveRecursiveArithmetic/ConciseIPRA/Lean/ConciseIPRA}{Pinned Lean sources}.
\bibitem{source} ProveIt contributors,
\emph{Source dispatch and literal-zero source elimination},
\code{Research/source_dispatch/} and \code{Lean/ConciseIPRA/NativeZeroSource.lean},
same pinned commit.
\href{https://github.com/VladimirReshetnikov/ProveIt/tree/1517d9d79ce33f631e0f82960c849427cbba23bf/Logic/PrimitiveRecursiveArithmetic/ConciseIPRA/Research/source_dispatch}{Research package}.
\bibitem{guarded} ProveIt contributors, \emph{A guarded core for iPRA},
\code{Research/guarded_core/}, same pinned commit. The unchanged Eq/$D$
comparison source is \code{reference/baseline_strict_goal.py} in that package.
\href{https://github.com/VladimirReshetnikov/ProveIt/tree/1517d9d79ce33f631e0f82960c849427cbba23bf/Logic/PrimitiveRecursiveArithmetic/ConciseIPRA/Research/guarded_core}{Pinned package}.
\bibitem{choice} ProveIt contributors,
\emph{Complementary introductions with one local assumption} and
\emph{Arithmetic strength and both ordinary interpretation composites},
\code{Research/certified_zero/reference/}, same pinned commit.
\bibitem{erasure} ProveIt contributors,
\emph{Internal scalar predecessor bootstrap},
\code{Research/certified_zero/reference/scalar_predecessor.tex},
same pinned commit; attributed predecessor construction reproduced above.
\bibitem{forward} ProveIt contributors,
\emph{Equality reflection and quantifier-free induction},
\code{Research/certified_zero/reference/equality_reflection.tex},
same pinned commit; attributed inherited reflection proof reproduced above.
\bibitem{progress} ProveIt contributors,
\emph{Primitive-recursion compilation},
\code{Research/certified_zero/reference/primitive_recursion.tex},
same pinned commit; attributed pairing and compilation proof reproduced above.
\end{thebibliography}
\end{document}
